\documentclass[11pt]{article}
\usepackage{amssymb,amsmath,verbatim,amsthm}
\newtheorem{Theorem}{Theorem}[section]
\newtheorem{Lemma}[Theorem]{Lemma}
\newtheorem{Corollary}[Theorem]{Corollary}
\newtheorem{Remark}[Theorem]{Remark}
\newtheorem{Example}[Theorem]{Example}
\newtheorem{Definition}[Theorem]{Definition}
\newtheorem{Construction}[Theorem]{Construction}
\newtheorem{Claim}[Theorem]{Claim}
\newtheoremstyle{noparens}
  {}{}
  {\itshape}{}
  {\bfseries}{.}
  { }
  {\thmname{#1}\thmnumber{ #2}\mdseries\thmnote{ #3}}
\theoremstyle{noparens}
\usepackage{enumerate}
\usepackage{color,xcolor}
\usepackage{indentfirst}
\setlength{\parindent}{2em}
%\indent ϣ����ʽ���Ķ����ı���
%%%%%%%%%%%%%%%%%%%%%%
%\newcommand{\Proof}{\noindent {\em Proof:\ }}
%\newcommand{\SZ}{\emptyset}
\newcommand{\ignore}[1]{}
\def\mod#1{\allowbreak \mkern18mu {\rm mod}\,#1}
\newcommand{\h}{\hphantom{{}\sp{*}}}
\newcommand{\vf}{\bigskip}
\newcommand{\tabsize}{\footnotesize}
\newcommand{\tabsizeoff}{\normalsize}

\def\widespacing{\relax}
\def\widespacingoff{\relax}
\textwidth 6.27in \textheight 8.99in \topmargin -0.5in
\parskip 6pt
\oddsidemargin=0.3in \evensidemargin=0.3in
\newcommand{\PSL}{{\operatorname{PSL}}}
\newcommand{\PGL}{{\operatorname{PGL}}}
\newcommand{\Fix}{{\operatorname{Fix}}}
\newcommand{\Sym}{{\operatorname{Sym}}}

\begin{document}

\renewcommand{\theequation}{\thesection.\arabic{equation}}

\title{\bf Steiner Quadruple Systems with Minimum Colorable Derived Designs: Constructions and Applications\footnote{The research was supported
 by National Natural Science Foundation of China (12171028, 12371326).}}
\author{
{\small Yuli  Tan,}  {\small  Junling  Zhou}\\
{\small School of Mathematics and Statistics}\\
{\small Beijing Jiaotong University}\\
{\small Beijing  100044, China}\\
{\small YL.Tan@bjtu.edu.cn}\\
{\small jlzhou@bjtu.edu.cn}\\
}
\date{ }
\maketitle

\begin{abstract}
An $r$-block-coloring, simply $r$-coloring, of a Steiner triple system $\mathrm{STS}(v)$ is a partition of the block set into $r$ color classes, each color class being a partial parallel class.
The chromatic index of $\mathrm{STS}(v)$, denoted by $\chi^{\prime}(v)$, is the smallest $r$ for which an $r$-coloring of an $\mathrm{STS}(v)$ exists.
A minimum colorable Steiner triple system $\mathrm{mcSTS}(v)$ is an $\mathrm{STS}(v)$ admitting a $\chi^{\prime} (v)$-coloring.
We generalize the notion of an $\mathrm{RDSQS}$ (a Steiner quadruple system $\mathrm{SQS}$ with resolvable derived designs) to $\mathrm{mcDSQS}$, % (an $\mathrm{SQS}$ with minimum colorable derived designs),
representing an $\mathrm{SQS}$ whose derived design at every point is minimum colorable.
This is motivated from an application in non-binary diameter perfect codes.
The purpose of this paper is to display a few recursive constructions to produce $\mathrm{mcDSQS}$s via Steiner systems $\mathrm{S}(3,K,v)$ with certain properties.
Among others, a  construction for $\mathrm{mcDSQS}$s is developed, which is also new even for $\mathrm{RDSQS}$s; special constructions concentrating only on $\mathrm{mcDSQS}(6n+2)$s are  demonstrated as well. As the main results, both
a new infinite family of $\mathrm{RDSQS}(6n+4)$s and the first infinite family of $\mathrm{mcDSQS}(6n+2)$s are constructed. To be specific, an $\mathrm{RDSQS}(2^{2m+1}+2)$
and an $\mathrm{mcDSQS}(2\cdot 9^{m}+2)$ are proved to exist, in which the former class gives rise to a new infinite family of large sets of Kirkman triple systems.
%A new infinite family of $\mathrm{mcDSQS}(2\cdot 9^{m}+2)$ is constructed.
%Then, we construct a new infinite family of $\mathrm{mcDSQS}(2\cdot 9^{n}+2)$ and a new infinite family of $\mathrm{RDSQS}(2^{2n+1}+2)$.
%We obtain the first infinite family of $\mathrm{mcDSQS}(2\cdot 9^{m}+2)$.
As applications, the smallest $q$ is  determined  such that a diameter perfect constant-weight $(n,\frac{1}{4}\tbinom{n}{3},6;4)_{q}$ code exists where
 $n \in\{ 2\cdot 9^{m}+2:m\geq 1\}\bigcup\{ 2^{2m+1}+2:m\geq 0\}$.
%As applications, $\frac{n}{2}+1$ is  shown to be the smallest $q$ such that a diameter perfect constant-weight $(n,\frac{1}{4}\tbinom{n}{3},6;4)_{q}$ code exists where $n \in\{ 2\cdot 9^{m}+2:m\geq 1\}$ and $\frac{n}{2}$ is  shown to be the smallest $q$ such that a diameter perfect constant-weight $(n,\frac{1}{4}\tbinom{n}{3},6;4)_{q}$ code exists where $n \in\{ 2^{2m+1}+2:m\geq 0\}$.
\medskip

\noindent {\bf Keywords}: Steiner quadruple system, minimum colorable, derived design, large set, Kirkman triple system, diameter perfect code
\medskip
\end{abstract}
%%%%%%%%%%%%%%%%%%%%%%%%%%%%%%%%%%%%%%%%%%%%%%%%%%%%%%%%%%%%%%%%%%%%%%%%%


\section{Introduction and preliminaries}

Let $X$ be a set of $v$ points and $K$ a set of positive integers.
Let $\mathcal{G}$ be a collection of subsets (called {\em{groups}} or {\em{holes}}) of $X$ which partition $X$.
A {\em{group divisible $t$-design}} of order $v$ with block sizes from $K$, denoted by $\mathrm{GDD}(t,K,v)$, is a triple $(X,\mathcal{G},\mathcal{B})$ with the following properties:\vspace{-0.2cm}
\begin{itemize}
  \item [(1)] $\mathcal{B}$ is a family of subsets of $X$ (called {\em{blocks}}), each having cardinality in $K$;\vspace{-0.2cm}
 \item [(2)] each block of $\mathcal{B}$ intersects any given group in at most one point;\vspace{-0.2cm}
 \item [(3)] each $t$-subset of $X$ from $t$ distinct groups is contained in exactly one block.
\end{itemize}
\vspace{-0.2cm}
The {\em{type}} of the $\mathrm{GDD}$ is defined to be the multiset $T=\{|G|: G \in \mathcal{G}\}$ of group sizes.
We also use $g_{1}^{a_{1}} g_{2}^{a_{2}}\ldots g_{r}^{a_{r}}$ to denote the type, which means that there are $a_{i}$ groups of size $g_{i}$, $1\leq i \leq r$.
A $\mathrm{GDD}(t,\{k\},v)$ is also denoted by $\mathrm{GDD}(t,k,v)$.
A $\mathrm{GDD}(t,k,kn)$ of type $n^{k}$ is known as a {\em{transversal $t$-design}} and denoted by $\mathrm{TD}(t,k,n)$.
If $(X,\mathcal{G},\mathcal{B})$ is a $\mathrm{GDD}(t,K,v)$ of type $1^{v}$, then we omit the notation of $\mathcal{G}$ and call $(X,\mathcal{B})$ a {\em{Steiner $t$-system}}, also denoted $\mathrm{S}(t,K,v)$. The notation $\mathrm{S}(t,k,v)$ is used for $K=\{k\}$.
An $\mathrm{S}(2,3,v)$ is a {\em{Steiner triple system}} $\mathrm{STS}(v)$ and an $\mathrm{S}(3,4,v)$ is a {\em{Steiner quadruple system}} $\mathrm{SQS}(v)$.
It is well-known that an $\mathrm{STS}(v)$ exists if and only if $v \equiv 1,3 \pmod{6}$ \cite{TSTS} and an $\mathrm{SQS}(v)$ exists if and only if $v \equiv 2,4 \pmod{6}$ \cite{HSQS}.

Let $(X,\mathcal{G},\mathcal{B})$ be a $\mathrm{GDD}$ and $Y\subseteq X$.
A {\em{partial parallel class}} ($\mathrm{PPC}$) of $Y$ is a collection of pairwise disjoint blocks contained in $Y$.
A {\em{parallel class}} ($\mathrm{PC}$) of $Y$ is a $\mathrm{PPC}$ that partitions $Y$.
We say that $(X, \mathcal{G}, \mathcal{B})$ is {\em {resolvable}}, denoted by $\mathrm{RGDD}$, if $\mathcal{B}$ can be partitioned into $\mathrm{PC}$s of $X$.
An $\mathrm{RGDD}(t,k,v)$ of type $1^{v}$ is a resolvable $\mathrm{S}(t,k,v)$, denoted by $\mathrm{RS}(t,k,v)$. An $\mathrm{RS}(2,3,v)$ is a {\em{Kirkman triple system}} $\mathrm{KTS}(v)$.
A $\mathrm{KTS}(v)$ exists if and only if $v\equiv  3 \pmod{6}$ \cite{KTS1,KTS2}. Two $\mathrm{KTS}(v)$s on the same set $X$ are said to be {\em{disjoint}} if they have no blocks in common. A {\em{large set of Kirkman triple system}} of order $v$, denoted by $\mathrm{LKTS}(v)$, is a partition of all triples of a $v$-element set $X$ into $v-2$ pairwise disjoint $\mathrm{KTS}(v)$s.
Although the problem of the existence of $\mathrm{LKTS}(v)$s has a long history, it remains pretty much open.

Let $(X,\mathcal{B})$ be an $\mathrm{S}(t,k,v)$. For any point $x\in X$, let $\mathcal{B}_{x} = \{B\setminus \{x\}:x\in B, B\in \mathcal{B}\}$. Then $(X\setminus\{x\},\mathcal{B}_{x})$ is an $\mathrm{S}(t-1,k-1,v-1)$, which is said to be the {\em{derived design}} of $(X,\mathcal{B})$ at the point $x$.
An $\mathrm{S}(t,k,v)$ with {\em resolvable derived designs}, abbreviated to $\mathrm{RDS}(t,k,v)$, refers to an $\mathrm{S}(t,k,v)$ whose derived design at every point is resolvable.
We use the notation $\mathrm{RDSQS}(v)$ to denote an $\mathrm{SQS}(v)$ with resolvable derived designs. $\mathrm{RDSQS}$s play an important role in recursive constructions of large sets of Kirkman triple systems (see \cite{Zhou-lkts,LEI1,Kang-2008}).

\begin{Theorem}[\label{RDS_1}{\cite{Zhou-lkts,Kang-2008}}]
If there exists an $\mathrm{RDSQS}(v + 1)$, then there exists an $\mathrm{LKTS}(3v)$.
\end{Theorem}

We know that the necessary condition for the existence of an $\mathrm{RDSQS}(v)$ is $v \equiv 4 \pmod{6}$ as a $\mathrm{KTS}(v)$ exists if and only if $v\equiv  3 \pmod{6}$.
 An $\mathrm{SQS}(v)$ with $v\equiv 2\pmod 6$ dose not have a resolvable derived design, but one may instead ask for a derived design with a partition of the blocks into minimum number of PPCs.
We introduce the concept of block-coloring.
An {\em{$r$-block-coloring}}, or simply {\em{$r$-coloring}}, of a design is a partition of the block set into $r$ color classes, each color class being a $\mathrm{PPC}$.
The collection of the $r$ $\mathrm{PPC}$s is called a {\em{resolution}}. The {\em{chromatic index}} of a design $\mathcal{D}$, denoted by $\chi^{\prime}(\mathcal{D})$, is the smallest $r$ for which $\mathcal{D}$ admits an $r$-coloring. The chromatic index of Steiner systems $\mathrm{S}(t,k,v)$, denoted by $\chi^{\prime}_{t,k}(v)$, is defined as
$$\chi^{\prime}_{t,k}(v) = \mathop{\min}_{\mathcal{D}\text{ is an }S(t,k,v)}\chi^{\prime}(\mathcal{D}).$$
For convenience, the notation $\chi^{\prime} (v)$ is used for $\chi^{\prime}_{2,3}(v)$.
If $\mathcal{D}$ is a resolvable $\mathrm{S}(t,k,v)$, then $\chi^{\prime}_{t,k}(\mathcal{D})$ equals  ${k\tbinom{v}{t}}/{v\tbinom{k}{t}}$. Therefore, $\chi^{\prime}(v)$ equals $(v-1)/2$ for $v\equiv  3 \pmod{6}$ because of the existence of a $\mathrm{KTS}(v)$.
The unique $\mathrm{STS}(7)$ has chromatic index $7$ as each pair of blocks has an intersection point.
There are two nonisomorphic $\mathrm{STS}(13)$s and each of them has chromatic index $8$ (see \cite{HTS14}). Generally, $\chi^{\prime}(v)$ equals $(v+1)/2$ if $v\equiv  1 \pmod{6}$ and $v\notin \{7,13\}$ (see \cite{HTS}). To sum up,
\begin{align}\label{E1}
\chi^{\prime}(v) =
\begin{cases}
\frac{v-1}{2}, & \text{ if } v\equiv  3 \pmod{6}, \\
\frac{v+1}{2}, & \text{ if } v\equiv  1 \pmod{6}\text{ and } v\notin \{7,13\},\\
7, & \text{ if } v=7,\\
8, & \text{ if } v=13.\\
\end{cases}
\end{align}
%\eqno(*)
%$$
%\begin{align*}
%\chi^{\prime}(v) = \begin{cases}
%\frac{v-1}{2}, & \text{ if } v\equiv  3 \pmod{6},\\
%\frac{v+1}{2}, & \text{ if } v\equiv  1 \pmod{6}\text{ and } v\notin \{7,13\},\\
%7, & \text{ if } v=7,\\
%8, & \text{ if } v=13.\\
%		   \end{cases}
%\end{align*}%\eqno(*)
%$$\chi(v) = \bigg\{\frac{v-1}{2}$$

If an $\mathrm{STS}(v)$ admits a $\chi^{\prime} (v)$-coloring, then it is said to be
{\em{minimum colorable}}, denoted by $\mathrm{mcSTS}(v)$.
If $v\equiv  3 \pmod{6}$, then an $\mathrm{mcSTS}(v)$ is a $\mathrm{KTS}(v)$.
An $\mathrm{SQS}(v)$ with {\em{minimum colorable derived designs}}, abbreviated to $\mathrm{mcDSQS}(v)$, refers to an $\mathrm{SQS}(v)$ whose derived design at every point is minimum colorable. Therefore, an $\mathrm{mcDSQS}(v)$ is the same as an $\mathrm{RDSQS}(v)$ if $v\equiv  4 \pmod{6}$. There is a large literature studying $\mathrm{RDSQS}$s \cite{zheng,Zhou-lkts,LEI1,yuanotripling}, while we only have one known $\mathrm{mcDSQS}(6n+2)$ with $n\geq 2$, that is, an $\mathrm{mcDSQS}(20)$ \cite{SHI}. In this work we introduce the concept of $\mathrm{mcDSQS}$s, not only as a generalization of $\mathrm{RDSQS}$s, but also due to their applications in diameter perfect codes.


Diameter perfect codes are a natural generalization of perfect codes. They are based on the code-anticode bound, a generalization of sphere-packing bound. Code-anticode bound was first introduced by Delsarte (see \cite{Delsarte}) and generalized by Ahlswede et al. (see \cite{DPD}).
Let $\mathcal{C}$ be a code in a metric space $\mathcal{V}$ with a transitive automorphism group, where the minimum distance between the codewords in $\mathcal{C}$ is $d$. Let $\mathcal{A}$ be a anticode of $\mathcal{V}$ such that the maximum distance (diameter) between elements in $\mathcal{A}$ is $D$, where $D< d$. Then
$$|\mathcal{C}|\cdot |\mathcal{A}|\leq |\mathcal{V}|.$$
If the inequality holds with equality, then $\mathcal{C}$ is said to be a {\em{$D$-diameter perfect code}} or {\em{diameter perfect distance-$d$ code}}.
Let $\mathcal{J}_{q}(n,w)$ denote the set of
all words of length $n$ and weight $w$, over $\mathbb{Z}_{q}=\{0,1,\ldots, q-1\}$, where the Hamming distance is used as the metric.
A code $\mathcal{C}\subseteq \mathcal{J}_{q}(n,w)$ is a {\em{$q$-ary constant-weight-$w$ code of length $n$}} and it is denoted by $(n,|\mathcal{C}|,d;w)_{q}$ code if $\mathcal{C}$ has minimum distance $d$.

In \cite{Etzion}, Etzion classified the known non-binary diameter perfect constant-weight codes into six families. Shi et al. \cite{SHI} generalized one of these families and investigated a family of such codes, where the supports of the codewords from a Steiner system $\mathrm{S}(t,w,n)$.
For $t\leq w$ and $2w-t\leq n$, define
$$\mathcal{A}_{q}(n,w,t)=\{(a_{1},\ldots,a_{t},\ldots,a_{n})\in \mathbb{Z}_{q}^{n}:\text{wt}((a_{1},\ldots,a_{t}))=t,\text{wt}((a_{t+1},\ldots,a_{n}))=w-t\}.$$
We can check that $\mathcal{A}_{q}(n,w,t)$ is a diameter-$(2w-t)$ anticode with ${{n-t}\choose{w-t}}(q-1)^{w}$ words in $\mathcal{J}_{q}(n,w)$.



\begin{Theorem}[\label{shi1}{\cite[Theorem 3]{SHI}}]
Let $t,w$ and $n$ be integers such that $0 <t \leq w < n$, $2w-t\leq n$.%\vspace{-0.4cm}
\begin{itemize}
  \item [$(1)$] An $(n,\tbinom{n}{t}/\tbinom{w}{t},2w-t+1;w)_{q}$ code is a diameter perfect constant-weight code which attains the code-anticode bound with respect to the anticode $\mathcal{A}_{q}(n,w,t)$;\vspace{-0.2cm}
  \item [$(2)$] If there exists a Steiner system $\mathrm{S}(t,w,n)$, then there exists an $(n,\tbinom{n}{t}/\tbinom{w}{t},2w-t+1;w)_{q}$ code with $q\geq \tbinom{n-1}{t-1}/\tbinom{w-1}{t-1}+1$.%\vspace{-0.2cm}
\end{itemize}
\end{Theorem}
Define $q_{0}^{\prime}(t,w,n)$ to be the smallest $q$ such that an $(n,\tbinom{n}{t}/\tbinom{w}{t},2w-t+1;w)_{q}$ code exists. Theorem \ref{shi1} (2) provides an upper bound of $q_{0}^{\prime}(t,w,n)$. When $t=3$ and $w=4$, the exact value of $q_{0}^{\prime}(t,w,n)$ is presented in terms of chromatic indexes of derived designs of an $\mathrm{SQS}(n)$.



\begin{Theorem}[\label{shi2}{\cite[Theorem 5]{SHI}}]
For $n\equiv 2,4 \pmod6$,
$$q_{0}^{\prime}(3,4,n)= \mathop{\min}_{S \text{ is } \mathrm{SQS}(n)}\mathop{\max}_{D\in S^{\prime}}\chi^{\prime}(D)+1,$$
where $S^{\prime}$ is the set of all $\mathrm{STS}(n-1)$s derived from $S$, the $\mathrm{SQS}(n)$.
\end{Theorem}
The following corollary is immediate.
\begin{Corollary}\label{shi-c1}
\begin{itemize}
 \item [$(1)$] For $n\equiv 2 \pmod6$ and $n\geq 20$, if there exists an $\mathrm{mcDSQS}(n)$, then $q_{0}^{\prime}(3,4,n)=\frac{n}{2}+1$.
 \item [$(2)$] For $n\equiv 4 \pmod6$, if there exists an $\mathrm{RDSQS}(n)$, then $q_{0}^{\prime}(3,4,n)=\frac{n}{2}$.
 \end{itemize}
\end{Corollary}


The rest of the paper is organized as follows. Section 2 is devoted to demonstrating recursive constructions for generic $\mathrm{mcDSQS}$s. In Subsection 2.1, we define a $\mathrm{gcDCQS}(g^{n}:s)$ to build an $\mathrm{mcDSQS}(gn+s)$ by filling in holes. Then we demonstrate several recursive constructions to produce $\mathrm{gcDCQS}(g^{n}:s)$ via Steiner systems $\mathrm{S}(3,K,v)$ with certain properties.
In Subsection 2.2, we concentrate on constructions of $\mathrm{mcDSQS}(6n+2)$s. We propose two new types of  colorings of the derived designs of a $\mathrm{CQS}(g^{n}:2)$ and then present related constructions.
In Section 3, we construct some necessary designs required in our recursive constructions. Then we establish the existence of a new infinite family of $\mathrm{RDSQS}(6n+4)$s and  an infinite family of $\mathrm{mcDSQS}(6n+2)$s. %give some new infinite families of $\mathrm{LKTS}$s.
%A new infinite family of $\mathrm{mcDSQS}(2\cdot 9^{m}+2)$ are constructed.
As  consequences, we show their applications in diameter perfect codes and large sets of Kirkman triple systems. In Section 4, we conclude the paper with some open questions.


\section{Recursive constructions}

A combinatorial structure named candelabra quadruple system $\mathrm{CQS}$ plays a central role in the recursive constructions of $\mathrm{SQS}$s. In order to produce
$\mathrm{RDSQS}$s, candelabra quadruple systems with resolvable derived designs $\mathrm{RDCQS}$s were introduced and intensively studied in the last decade. %\cite{Zhou-lkts,LEI1}.
In this section, we introduce several coloring types of the derived designs of a $\mathrm{CQS}$ as a generalization of an $\mathrm{RDCQS}$. Then we display a number of recursive constructions to produce $\mathrm{mcDSQS}$s. For later use, we introduce or recall some definitions.


An {\em{incomplete Steiner triple system}} on a $v$-set $X$ with an $h$-subset $H$ as a hole, denoted by $\mathrm{STS}(v,h)$, is a triple $(X,H,\mathcal{B})$, where $\mathcal{B}$ is a block set of a $\mathrm{GDD}(2,3,v)$ of type $1^{v-h}h^{1}$ with $H$ as the long group.
Further let $v\equiv 1,3\pmod 6$.
If $\mathcal{B}$ can be partitioned into $r:=\chi^{\prime}(v)$ $\mathrm{PPC}$s of $X$, in which there are $s:=\chi^{\prime}(v)-\frac{v-h}{2}$ $\mathrm{PPC}$s of $X\setminus H$,
then we say that $(X,H,\mathcal{B})$ is {\em{good colorable}}, denoted by $\mathrm{gcSTS}(v,h)$.
It then follows from (\ref{E1}) that
\begin{align*}
(r,s) = \begin{cases}
(\frac{v-1}{2},\frac{h-1}{2}), & \text{ if } v\equiv  3 \pmod{6},\\
(\frac{v+1}{2},\frac{h+1}{2}), & \text{ if } v\equiv 1 \pmod{6}, v\neq 7,13.\\
		   \end{cases}
\end{align*}
It is not difficult to check that, when $v\equiv 3\pmod 6$, a $\mathrm{gcSTS}(v,h)$ $(X,H,\mathcal{B})$ becomes an {\em incomplete Kirkman triple system }$\mathrm{KTS}(v,h)$, whose blocks are partitioned into $\frac{v-h}{2}$ $\mathrm{PC}$s of $X$ and $\frac{h-1}{2}$ $\mathrm{PC}$s of $X\setminus H$.
%Note that an $\mathrm{STS}(v,h)$ has $b$ blocks, where $b=\frac{v(v-1)-h(h-1)}{6}=\frac{v-h}{2}\cdot\frac{v}{3}+\frac{h-1}{2}\cdot\frac{v-h}{3}$. Hence a $\mathrm{gcSTS}(v,h)$ cannot exist if $(v,h)\equiv(3,1) \pmod 6$ and $h\neq 1$. So when we talk about a $\mathrm{gcSTS}(v,h)$, we always assume that
%\begin{align}
%\begin{array}{c}\label{aa}
 %   		v,h\equiv 1,3\hspace{-3mm}\pmod 6\text{ but } (v,h)\not\equiv(3,1) \hspace{-3mm}\pmod 6 \text { if } h\neq 1,\\
 %           v\geq 19 \text { if } v\equiv 1 \hspace{-3mm}\pmod 6.
  %  	\end{array}
%\end{align}



An {\em{incomplete  Steiner quadruple system}} on a $v$-set $X$ with an $h$-subset $H$ as a hole, denoted by $\mathrm{SQS}(v,h)$, is a triple $(X,H,\mathcal{B})$, where $\mathcal{B}$ is a set of quadruples of $X$ (blocks), such that every triple $T\subseteq X$ with $T\not\subseteq H$ is contained in a unique block and no triple of $H$ is contained in any block.
%where $\mathcal{B}$ is a block set of a $\mathrm{GDD}(3,4,v)$ of type $1^{v-h}h^{1}$ with $H$ as the long group.
We can also define the derived designs of an $\mathrm{SQS}(v,h)$.
For $x \in X$, put $\mathcal{B}_x=\{B \setminus\{x\}: x \in B, B \in \mathcal{B}\}$.
It is immediate that each $(X \setminus\{x\},\mathcal{B}_x)$ with $x\in X\setminus H$ forms an $\mathrm{STS}(v-1)$ and each $(X \setminus \{x\}, H\setminus\{x\}, \mathcal{B}_x)$ with $x \in H$ forms an $\mathrm{STS}(v-1,h-1)$.
An $\mathrm{SQS}(v,h)$ with {good colorable derived designs}, denoted by $\mathrm{gcDSQS}(v,h)$, if all derived $\mathrm{STS}(v-1)$s are minimum colorable and all derived $\mathrm{STS}(v-1,h-1)$s are good colorable.


A {\em{Kirkman frame}} $\mathrm{F}(2,3,v\{m\})$ (as in \cite{frame}), briefly by $\mathrm{KF}(m^{v})$, is a $\mathrm{GDD}(2,3,mv)$ of type $m^{v}$ $(X,\mathcal{G},\mathcal{B})$ such that the block set $\mathcal{B}$ can be partitioned into subsets $P_{r}$, $r\in R$, each $P_{r}$ being a $\mathrm{PC}$ of $X\setminus G$ for some group $G\in \mathcal{G}$. A $\mathrm{KF}(m^{v})$ exists if and only if $m\equiv 0 \pmod 2$, $v\geq 4$ and $m(v-1)\equiv 0 \pmod 3$ (see \cite{KF}).

Let $(X, \mathcal{G}, \mathcal{A})$ be a $\mathrm{GDD}(3,4, mv)$ of type $m^v$. For any $x \in X$, let $\mathcal{A}_x=\{A \backslash\{x\}: x \in A, A \in \mathcal{A}\}$. Then the derived design $(X \backslash G, \mathcal{G} \backslash\{G\}, \mathcal{A}_x)$ is a $\mathrm{GDD}(2,3, m(v-1))$ of type $m^{v-1}$.
Here we examine two types of coloring of the derived designs. We use $\mathrm{RDGDD}(3,4, v\{m\})$ to denote a $\mathrm{GDD}(3,4, mv)$ of type $m^{v}$ with resolvable derived designs, meaning that its derived design at every point is a resolvable $\mathrm{GDD}(2,3,m(v-1))$.
An $\mathrm{RDTD}(3,4,m)$ denotes an $\mathrm{RDGDD}(3,4,4\{m\})$.
Parallelly, we use $\mathrm{FDGDD}(3,4,v\{m\})$ to denote a $\mathrm{GDD}(3,4,mv)$ of type $m^{v}$ whose derived design at every point is a $\mathrm{KF}(m^{v-1})$.

Let $X$ be a $v$-element set. Suppose that $\mathcal{G}$, $\mathcal{B}_{i}$ $(1\leq i \leq s)$ and $\mathcal{T}$ are collections of some subsets of $X$ and $\mathcal{G}$ is a partition of $X$.
Let $K_{i}$ $(1\leq i \leq s)$ and $K_{T}$ be the block size sets of $\mathcal{B}_{i}$ and $\mathcal{T}$ respectively.
An {\em{$s$-fan design}}, denoted by $s$-$\mathrm{FG}(3,(K_1, \ldots, K_s, K_T), v)$, is an $(s+3)$-tuple $(X,\mathcal{G},\mathcal{B}_{1},\ldots,\mathcal{B}_{s},\mathcal{T})$ such that each $(X,\mathcal{G}\cup \mathcal{B}_{i})$ with $1\leq i \leq s$ is a Steiner $2$-system and $(X,\mathcal{G}\bigcup (\bigcup_{i=1}^{s}\mathcal{B}_{i})\bigcup \mathcal{T})$ is a Steiner $3$-system. %$\mathrm{S}(3,(\bigcup_{i=1}^{s}K_{i})\bigcup K_{T},v)$.
The type of an $s$-$\mathrm{FG}$ is the multiset $\{|G|: G \in \mathcal{G}\}$. Note that an $s$-$\mathrm{FG}(3,(K_1, \ldots, K_s, K_T), v)$
of type $1^v$ is an $\mathrm{S}(3,(\bigcup_{i=1}^{s}K_{i})\bigcup K_{T},v)$ having $s$ subdesigns $\mathrm{S}(2,K_{i},v)$. For such a case we simply say that $(X,\mathcal{B}_{1},\ldots,\mathcal{B}_{s},\mathcal{T})$ is an $s$-fan $\mathrm{S}(3,(K_{1},\ldots, K_{s},$ $K_{T}),v)$.






%Let $K$ be a set of positive integers.
 For a $v$-set $X$, take an $s$-subset $S\subseteq X$ and call it a {\em{stem}}. Let $\mathcal{G}$ be a collection of subsets, called groups, of $X\setminus S$ which partition $X\setminus S$.
A {\em{candelabra $t$-system}} (as in \cite{Cande}) of order $v$ with block sizes from a set $K$, briefly by $\mathrm{CS}(t,K,v)$, is a quadruple $(X, S, \mathcal{G}, \mathcal{A})$ satisfying that
\vspace{-0.2cm}
\begin{itemize}
  \item [(1)] $\mathcal{A}$ is a family of subsets of $X$, called blocks, each of cardinality from $K$; \vspace{-0.2cm}
  \item [(2)] each $t$-subset  of $X$ except those contained in $S\cup G$ for some $G\in \mathcal{G}$ is contained in exactly one block (while no $t$-subset  in $S\cup G$ is contained in any block).
\end{itemize}
\vspace{-0.2cm}
By the type of a $\mathrm{CS}(t,K,v)$ we mean the list $(\{|G|: G \in \mathcal{G}\}: s)$. Exponential notation is also applied to denote the type of $\mathcal{G}$. A candelabra system with $t=3$ and $k=4$ is a {\em{candelabra quadruple system}} and denoted by $\mathrm{CQS}$. %A candelabra system with $t=3$ and $k=3$ is denoted by $\mathrm{CS}$.









It is a standard construction to build an $\mathrm{SQS}(gn+s)$ from a $\mathrm{CQS}(g^{n}:s)$ by filling in holes. To produce $\mathrm{mcDSQS}$s, we need to define desirable
 colorings of the derived designs of  $\mathrm{CQS}$s. The rest of this section is divided into two subsections. In the first subsection, we generalize  known
  definitions and constructions related with $\mathrm{RDSQS}$s to develop recursive constructions applicable for generic $\mathrm{mcDSQS}$s. %We define a $\mathrm{CQS}$ with good colorable derived designs and then demonstrate three constructions via Steiner systems $\mathrm{S}(3,k,v)$ with certain properties.
 In the second subsection, we turn our attention to $\mathrm{mcDSQS}$s only of orders $6n+2$. We introduce two new types of colorings of derived designs of a $\mathrm{CQS}(g^{n}:2)$ and display effective recursive constructions.



\subsection{Generic constructions}




Let $(X, S, \mathcal{G}, \mathcal{A})$ be a $\mathrm{CQS}(g^n: s)$.
For $x \in X$, put $\mathcal{A}_x=\{A \setminus\{x\}: x \in A, A \in \mathcal{A}\}$. It can be checked that each derived design $(X \setminus\{x\},(G \cup S) \setminus\{x\}, \mathcal{A}_x)$ with $x \in G$ and $G \in \mathcal{G}$ forms an $\mathrm{STS}(gn+s-1, g+s-1)$; and each derived design $(X \setminus S, \mathcal{G}, \mathcal{A}_x)$ with $x \in S$ forms a $\mathrm{GDD}(2,3, gn)$ of type $g^n$.
A $\mathrm{CQS}(g^n: s)$ is said to have {\em good colorable derived designs}, denoted by $\mathrm{gcDCQS}(g^n: s)$, if the collection $\{\mathcal{A}_x:x \in X\}$ of its derived designs has the following properties:\vspace{-0.2cm}
\begin{itemize}
  \item [(1)] for $x\in G$ and $G\in \mathcal{G}$, each $(X\setminus \{x\}, (G\cup S)\setminus \{x\},\mathcal{A}_{x})$ is a $\mathrm{gcSTS}(gn+s-1,g+s-1)$;\vspace{-0.2cm}
  \item [(2)] for $x\in S$, each $(X\setminus S, \mathcal{G}, \mathcal{A}_{x})$ is a $\mathrm{KF}(g^{n})$. \vspace{-0.2cm}
\end{itemize}

Whenever $gn+s\equiv 4\pmod 6$ and $s\ne 1$, a $\mathrm{gcDCQS}(g^{n}:s)$ becomes an $\mathrm{RDCQS}(g^{n}:s)$, the notation used in \cite{Zhou-lkts}.
A simple but important use of $\mathrm{gcDCQS}$s is the filling construction to build $\mathrm{mcDSQS}$s.
%There are some necessary conditions for the existence of a $\mathrm{gcDCQS}(g^{n}:s)$. If $g\equiv 0\pmod 2$ and $gn+s \not\equiv g+s \pmod 6$, then $g(n-1)\not\equiv 0 \pmod 3$ and thus there is no $\mathrm{gcDCQS}(g^{n} : s)$ if $s \geq1$, since a $\mathrm{KF}(g^{n})$ does not exist. Combining with (\ref{aa}), when we use a $\mathrm{gcDCQS}(g^{n}:s)$, we have the following:
%\begin{eqnarray}
%\begin{array}{c}\label{bb}
%    		gn+s\equiv g+s \equiv 2,4 \hspace{-3mm}\pmod 6, \text{ when } s\geq 2,\\
%            g, gn \equiv 2,4 \hspace{-3mm}\pmod 6 \text{ but } gn\not\equiv 4 \hspace{-3mm}\pmod 6 \text{ if } g \equiv 2\hspace{-3mm}\pmod6, \text{ when } s=0.\\
%\end{array}
%\end{eqnarray}
{\Construction \label{mcdcqs-mcdsqs}
If there exist a $\mathrm{gcDCQS}(g^n: s)$, an $\mathrm{mcDSQS}(g+s)$
 and a $\mathrm{gcDSQS}(g+s,s)$, then there exists an $\mathrm{mcDSQS}(gn+s)$.}

\proof %Note the conditions in (\ref{bb}). %%������û������
See \cite{Zhou-lkts} for the case $gn+s \equiv 4 \pmod 6$.
We only need to handle the case $gn+s \equiv  2 \pmod 6$.
Suppose $(X,S,\mathcal{G},\mathcal{A})$ is the given $\mathrm{gcDCQS}(g^n: s)$.
For $x \in X$, let $\mathcal{A}_x=\{A \setminus\{x\}: x \in A, A \in \mathcal{A}\}$.
For $x\in G$ and $G\in \mathcal{G}$, each $(X\setminus \{x\}, (G\cup S)\setminus \{x\},\mathcal{A}_{x})$ is a $\mathrm{gcSTS}(gn+s-1,g+s-1)$; its resolution is $\{P_{x}^{j}:1 \leq j \leq \frac{gn+s}{2}\}$, where each
$P_{x}^{j}$ is a $\mathrm{PPC}$ of $X\setminus \{x\}$;
particularly, each $P_{x}^{j}$ with $1 \leq j \leq \frac {g+s}{2}$ is a $\mathrm{PPC}$ of $X\setminus (G\cup S)$.
For $x\in S$ , each $(X\setminus S, \mathcal{G}, \mathcal{A}_{x})$ is a $\mathrm{KF}(g^{n})$; its resolution is $\{P_{x}^{j}(G):G\in \mathcal{G}, 1\leq j\leq \frac{g}{2}\}$, where each $P_{x}^{j}(G)$ is a $\mathrm{PC}$ of $X\setminus (G\cup S)$.

Let $G_{0}\in \mathcal{G}$. We can construct by assumption an $\mathrm{mcDSQS}(g+s)$ on $G_{0}\cup S$ with block set $\mathcal{B}(G_{0})$. So we have $g+s$ derived $\mathrm{mcSTS}(g+s-1)$s, say $((G_{0}\cup S) \setminus\{x\}, \mathcal{B}_{x}(G_{0}))$ with $x\in G_{0}\cup S$.
Further, each $\mathcal{B}_{x}(G_{0})$ has a resolution $\{Q_{x}^{j}(G_{0}):1\leq j \leq \frac {g+s}{2}\}$ into $\mathrm{PPC}$s of $(G_{0}\cup S)\setminus \{x\}$.


For each $G\in \mathcal{G}$ with $G\neq G_{0}$, we construct a $\mathrm{gcDSQS}(g+s,s)$ on $G\cup S$ with $S$ as a hole and with block set $\mathcal{B}(G)$. So we have $g$ derived $\mathrm{mcSTS}(g+s-1)$s $((G\cup S) \setminus\{x\}, \mathcal{B}_{x}(G))$ with $x\in G$ and $s$ derived $\mathrm{gcSTS}(g+s-1,s-1)$s $((G\cup S) \setminus\{x\}, S\setminus \{x\},\mathcal{B}_{x}(G))$ with $x\in S$.
Further, each $\mathcal{B}_{x}(G)$ with $x\in G$ has a resolution $\{Q_{x}^{j}(G):1\leq j \leq \frac {g+s}{2}\}$ into $\mathrm{PPC}$s of $(G\cup S)\setminus \{x\}$.
Each $\mathcal{B}_{x}(G)$ with $x\in S$ has a resolution $\{Q_{x}^{j}(G):1\leq j \leq \frac{g+s}{2}\}$ into $\mathrm{PPC}$s of $(G\cup S)\setminus\{x\}$; particularly, each $Q_{x}^{j}(G)$ with $\frac{g}{2}+1\leq i \leq \frac{g+s}{2}$ is a $\mathrm{PPC}$ of $G$.
Define
$$\mathcal{F} = \mathcal{A}\bigcup\left(\bigcup_{G\in\mathcal{G}}\mathcal{B}(G)\right).$$
It is straightforward that $(X,\mathcal{F})$ forms an $\mathrm{SQS}(gn+s)$.
For $x\in G$ and $G\in \mathcal{G}$, the derived design at $x$ has the block set
$\mathcal{A}_{x}\bigcup \mathcal{B}_{x}(G)$ and a resolution
$$\bigg\{P_{x}^{j}\cup Q_{x}^{j}(G):1\leq j \leq\frac {g+s}{2}\bigg\}\bigcup\bigg\{P_{x}^{j}:\frac {g+s}{2}+1 \leq j \leq \frac {gn+s}{2}\bigg\}.$$
For each $x\in S$, the derived design at $x$ has the block set
$\mathcal{A}_{x}\bigcup (\bigcup_{G\in \mathcal{G}}\mathcal{B}_{x}(G))$ and a resolution
$$\bigg\{P_{x}^{j}(G)\cup Q_{x}^{j}(G):G\in \mathcal{G},1\leq j \leq\frac {g}{2}\bigg\}\bigcup\bigg\{\bigcup_{G\in \mathcal{G}} Q_{x}^{j}(G):\frac {g}{2}+1 \leq j \leq \frac {g+s}{2}\bigg\}.$$
Hence, we end up with an $\mathrm{mcDSQS}(gn+s)$.
\qed










{\Construction \label{C3} Let $mn+s\equiv mk+s\equiv 2,4\pmod 6$. Suppose that there exists an $\mathrm{RDS}(3,k+1,n+1)$. If there exist a $\mathrm{gcDCQS}(m^{k}: s)$ and an $\mathrm{RDGDD}(3,4, (k+1)\{m\})$, then there exists a $\mathrm{gcDCQS}(m^n: s)$.}
\proof See \cite[Lemma 5.5]{Zhou-lkts} for the case $mn+s \equiv mk+s\equiv 4\pmod 6$. Next, we prove the case  $mn+s \equiv mk+s\equiv 2\pmod 6$. Suppose $(I_n \cup\{\infty\}, \mathcal{B})$ is the given $\mathrm{RDS}(3, k+1, n+1)$. Let $\mathcal{B}_\infty=\{B \backslash\{\infty\}: \infty\in B, B \in \mathcal{B}\}$ and $\mathcal{T}=\{B: \infty \notin B, B \in \mathcal{B}\}$. Thus $(I_n, \mathcal{B}_{\infty}, \mathcal{T})$ forms a $1$-fan $\mathrm{S}(3,(\{k\}, \{k+1\}), n)$, where $(I_{n},\mathcal{B}_{\infty})$ is an $\mathrm{S}(2,k,n)$. Let $S=\{\alpha_1, \alpha_2,\ldots, \alpha_s\}, Y=(I_n \times I_m) \cup S$, $\mathcal{H}=\{\{x\} \times I_m: x \in I_n\}$ and $\mathcal{H}_{B} = \{\{x\}\times I_{m}:x\in B\}$, where $B\subseteq I_{n}$.

For any $B\in \mathcal{B}_{\infty}$, we can construct by assumption a $\mathrm{gcDCQS}(m^{k}: s)$ $((B\times I_{m})\cup S,S,\mathcal{H}_{B},\mathcal{A}_B)$; its derived designs are $\mathcal{A}_{B}(x,i)$ with $(x,i)\in B\times I_{m}$ and $\mathcal{A}_{B}(\alpha)$ with $\alpha\in S$.
Moreover, the following hold:\vspace{-0.2cm}
\begin{itemize}
  \item [$\bullet$] each $\mathcal{A}_{B}(x, i)$, where $(x, i) \in B \times I_m$ and $B \in \mathcal{B}_{\infty}$, is the block set of a $\mathrm{gcSTS}(mk+s-1, m+s-1)$ on $((B \times I_m) \setminus\{(x, i)\}) \cup S$ with a hole $(\{x\} \times(I_m \setminus\{i\})) \cup S$;
       its resolution is $\{\mathcal{A}_{B}^{j}(x, i):1\leq j \leq \frac{mk+s}{2}\}$, where each $\mathcal{A}_{B}^{j}(x, i)$ is a $\mathrm{PPC}$ of $((B \times I_m) \setminus\{(x, i)\}) \cup S$; particularly, those $\mathcal{A}_{B}^{j}(x, i)$ with $\frac{m(k-1)}{2} +1 \leq j\leq \frac{mk+s}{2}$ are $\mathrm{PPC}$s of  $(B \setminus\{x\})\times I_{m}$;\vspace{-0.2cm}
  \item [$\bullet$] each $(B\times I_{m},\mathcal{H}_{B},\mathcal{A}_B(\alpha))$, where $\alpha\in S$, is a $\mathrm{KF}(m^k)$; its resolution is $\{\mathcal{A}_{B}^{y,j}(\alpha):y\in B, 1\leq j\leq \frac{m}{2}\}$, where each $\mathcal{A}_{B}^{y,j}(\alpha)$ is a $\mathrm{PC}$ of $(B \setminus\{y\})\times I_{m}$.\vspace{-0.2cm}
\end{itemize}

For any $B\in \mathcal{T}$, we can construct by assumption an $\mathrm{RDGDD}(3,4,(k+1)\{m\})$ $(B\times I_{m},\mathcal{H}_{B},\mathcal{D}_B)$; its derived designs are $\mathcal{D}_{B}(x,i)$ with $(x,i)\in B\times I_{m}$.
Moreover,
each $\mathcal{D}_B(x, i)$ forms the block set of an $\mathrm{RGDD}(2,3, mk)$ with group set $\mathcal{H}_{B\setminus\{x\}}$; its resolution is $\{\mathcal{D}_B^j(x, i): 1 \leqslant j \leqslant \frac{m(k-1)}{2}\}$, each $\mathcal{D}_B^j(x, i)$ being a $\mathrm{PC}$ of $(B \setminus\{x\}) \times I_m$.
Define
$$\mathcal{F} = \left(\bigcup_{B\in \mathcal{B}_{\infty}}\mathcal{A}_{B}\right)\bigcup\left(\bigcup_{B\in \mathcal{T}}\mathcal{D}_{B}\right).$$
Then it is routine to check that $(Y,S,\mathcal{H},\mathcal{F})$ forms a $\mathrm{CQS}(m^{n}:s)$. Next we verify that $(Y,S,\mathcal{H},\mathcal{F})$ is a $\mathrm{gcDCQS}(m^n: s)$.
For $(x,i)\in I_{n}\times I_{m}$, the derived design at $(x,i)$ has the block set
$$\mathcal{F}(x,i) = \left(\bigcup_{x\in B,B\in \mathcal{B}_{\infty}}\mathcal{A}_{B}(x,i)\right)\bigcup\left(\bigcup_{x\in B,B\in \mathcal{T}}\mathcal{D}_{B}(x,i)\right).$$
For $\alpha\in S$, the derived design at $\alpha$ has the block set
$$\mathcal{F}(\alpha)=\bigcup_{B\in \mathcal{B}_{\infty}}\mathcal{A}_{B}(\alpha).$$


For each $x \in I_n\cup \{\infty\}$, let $\mathcal{B}_x=\{B \backslash\{x\}: x \in B, B \in \mathcal{B}\}$. Since $(I_{n}\cup\{\infty\},\mathcal{B})$ is an $\mathrm{RDS}(3,k+1,n+1)$, each $((I_n \cup\{\infty\}) \setminus\{x\}, \mathcal{B}_x)$ forms an $\mathrm{RS}(2, k, n)$. Let $\Gamma_x=\{\mathcal{B}_x(h): 1 \leqslant h \leqslant p\}$, where $p=(n-1) /(k-1)$, be the resolution of $\mathcal{B}_{x}$.
For any $x \in I_n$, $1 \leqslant h \leqslant p$, there is a unique block $M_{x}(h)\in \mathcal{B}_{x}(h)$ containing $\infty$. Therefore, for any $B \in \mathcal{B}_x(h) \setminus\{M_{x}(h)\}$, we must have $B \cup\{x\} \in \mathcal{T}$. Denote $M_{\infty}(h)=(M_{x}(h)\cup\{x\}) \backslash\{\infty\}$. Then $\{M_{\infty}(h) \setminus\{x\}: 1 \leqslant h \leqslant p\}=\{M_{x}(h) \backslash\{\infty\}: 1 \leqslant h \leqslant p\}$ forms a partition of $I_n \backslash\{x\}$.


(1) Consider the coloring of the derived design $\mathcal{F}(x, i),(x, i) \in I_n \times I_m$, which forms the block set of an $\mathrm{STS}(mn+s-1, m+s-1)$ on $Y \setminus\{(x, i)\}$ with a hole $((\{x\} \times I_m)\cup S) \setminus\{(x,i)\}$.

For $1\leq j \leq \frac{m(k-1)}{2}$ and $1\leq h\leq p$, let
$$\mathcal{F}_{h}^{j}(x,i) = \mathcal{A}_{M_{\infty}(h)}^{j}(x,i) \bigcup \left( \bigcup _{B\in \mathcal{B}_{x}(h)\setminus \{M_{x}(h)\}}\mathcal{D}_{B\cup\{x\}}^{j}(x,i)\right).$$

For $\frac{m(k-1)}{2}+1\leq j \leq  \frac{mk+s}{2}$, let
$$\mathcal{F}^{j}(x,i) = \bigcup_{h=1}^{p}\mathcal{A}_{M_{\infty}(h)}^{j}(x,i).$$
It is immediate that the collection $\{\mathcal{F}_{h}^{j}(x,i):1\leq j \leq \frac{m(k-1)}{2},1\leq h\leq p\}\bigcup \{\mathcal{F}^{j}(x,i):\frac{m(k-1)}{2}+1\leq j \leq  \frac{mk+s}{2}\}$ forms a resolution of $\mathcal{F}(x,i)$ into $\mathrm{PPC}$s of $Y\setminus \{(x,i)\}$; in particular, each $\mathcal{F}^{j}(x,i)$ with $\frac{m(k-1)}{2}+1\leq j \leq  \frac{mk+s}{2}$ is a $\mathrm{PPC}$ of $(I_{n}\setminus \{x\} )\times I_m$. As a result, $(Y\setminus\{(x,i)\},((\{x\} \times I_m)\cup S) \setminus\{(x,i)\},\mathcal{F}(x, i))$ is good colorable.



(2) Consider the coloring of the derived design $\mathcal{F}(\alpha)$, $\alpha\in S$, which forms the block set of a $\mathrm{GDD}(2,3,mn)$ of type $m^{n}$ on $I_{n}\times I_{m}$ with group set $\mathcal{H}$.

For $y\in I_{n}$ and $1\leq j\leq \frac{m}{2}$, let
$$\mathcal{F}^{y,j}(\alpha) =  \bigcup_{y\in B,B\in \mathcal{B}_{\infty}}\mathcal{A}_{B}^{y,j}(\alpha).$$
We can check that $\mathcal{F}(\alpha)$ is the block set of a $\mathrm{KF}(m^{n})$ with a resolution $\{\mathcal{F}^{y,j}(\alpha):y\in I_{n},1\leq j\leq \frac{m}{2}\}$,
where each $\mathcal{F}^{y,j}(\alpha)$ is a $\mathrm{PC}$ of $
(I_{n}\setminus\{y\})\times I_{m}$.
Hence, we end up with a $\mathrm{gcDCQS}(m^n: s)$.
\qed

{\Construction \label{C2} Suppose that there exists a $1$-$\mathrm{FG}(3,(K_{1},K),gn)$ of type $g^{n}$.
Let $m,s$ be integers such that $mgn+s\equiv mk_{1}+s\equiv 2,4 \pmod 6$ for any $k_{1}\in K_{1}$.
 If there exist a $\mathrm{gcDCQS}(m^{k_{1}}: s)$ for any $k_{1}\in K_{1}$ and an $\mathrm{FDGDD}(3,4, k\{m\})$ for any $k\in K$, then there exists a $\mathrm{gcDCQS}((mg)^n: s)$.}
\proof
See \cite[Lemma 5.14]{Zhou-lkts} for the case $mgn+s\equiv mk_{1}+s\equiv 4\pmod 6$. Next, we only prove the case  $mgn+s\equiv mk_{1}+s\equiv 2\pmod 6$. %the case $mgn+s\equiv mg+s\equiv 2\pmod 6$ and $m+s\equiv 4\pmod 6$ is similar.
Suppose $(X,\mathcal{G},\mathcal{B},\mathcal{T})$ is the given $1$-$\mathrm{FG}(3,(K_{1},K),gn)$ of type $g^{n}$. Let $S=\{\alpha_{1},\alpha_{2},\ldots, \alpha_{s}\}$, $Y=(X \times I_m) \cup S$, $\mathcal{H} = \{G\times I_{m}:G\in \mathcal{G}\}$ and $\mathcal{H}_{B} = \{\{x\}\times I_{m}:x\in B\}$, where $B\subseteq X$.

For any $B\in \mathcal{B}$, construct a $\mathrm{gcDCQS}(m^{|B|}: s)$ $((B\times I_{m})\cup S,S, \mathcal{H}_{B}, \mathcal{A}_{B})$; its derived designs are
$\mathcal{A}_{B}(x,i)$ with $(x,i)\in B\times I_{m}$ and $\mathcal{A}_{B}(\alpha)$ with $\alpha\in S$. Moreover, we have the following:\vspace{-0.2cm}
\begin{itemize}
  \item [$\bullet$] each $\mathcal{A}_{B}(x, i)$, where $(x, i) \in B \times I_m$, is the block set of a $\mathrm{gcSTS}(m|B|+s-1, m+s-1)$ on $((B \times I_m) \setminus\{(x, i)\}) \cup S$ with a hole $(\{x\} \times(I_m \setminus\{i\})) \cup S$;
       its resolution is $\{\mathcal{A}_{B}^{y,j}(x, i):y\neq x, y\in B,1\leq j \leq \frac{m}{2}\}\cup\{\mathcal{A}_{B}^{j}(x, i):1 \leq j\leq \frac{m+s}{2}\}$, where each $\mathcal{A}_{B}^{y,j}(x, i)$ is a $\mathrm{PPC}$ of $((B \times I_m) \setminus\{(x, i)\}) \cup S$, and those $\mathcal{A}_{B}^{j}(x, i)$ are $\mathrm{PPC}$s of  $(B \setminus\{x\})\times I_{m}$;\vspace{-0.2cm}

  \item [$\bullet$] each $(B\times I_{m},\mathcal{H}_{B},\mathcal{A}_B(\alpha))$, where $\alpha\in S$, is a $\mathrm{KF}(m^{|B|})$; its resolution is $\{\mathcal{A}_{B}^{y,j}(\alpha):y\in B, 1\leq j\leq \frac{m}{2}\}$, where each $\mathcal{A}_{B}^{y,j}(\alpha)$ is a $\mathrm{PC}$ of $(B \setminus\{y\})\times I_{m}$.\vspace{-0.2cm}
\end{itemize}

For any $B\in \mathcal{T}$, we can construct by assumption an $\mathrm{FDGDD}(3,4,|B|\{m\})$ $(B\times I_{m},\mathcal{H}_{B},\mathcal{D}_B)$; its derived designs are $\mathcal{D}_{B}(x,i)$ with $(x,i)\in B\times I_{m}$.
Moreover, each $((B\setminus\{x\})\times I_{m},\mathcal{H}_{B\setminus\{x\}},$ $\mathcal{D}_{B}(x,i))$ is a $\mathrm{KF}(m^{|B|-1})$; its resolution is $\{\mathcal{D}_{B}^{y,j}(x,i):y\in B\setminus\{x\}, 1\leq j\leq \frac{m}{2}\}$, where each $\mathcal{D}_{B}^{y,j}(x,i)$ is a $\mathrm{PC}$ of $(B \setminus\{x,y\})\times I_{m}$.
Define

$$\mathcal{F} = \left(\bigcup_{B\in \mathcal{B}}\mathcal{A}_{B}\right)\bigcup\left(\bigcup_{B\in \mathcal{T}}\mathcal{D}_{B}\right).$$
Then it is routine to check that $(Y,S,\mathcal{H},\mathcal{F})$ forms a $\mathrm{CQS}((mg)^{n}:s)$. Next we verify that $(Y,S,\mathcal{H},\mathcal{F})$ is a $\mathrm{gcDCQS}((mg)^n: s)$.
For $(x,i)\in X\times I_{m}$, the derived design at $(x,i)$ has the block set
$$\mathcal{F}(x,i) = \left(\bigcup_{x\in B,B\in \mathcal{B}}\mathcal{A}_{B}(x,i)\right)\bigcup\left(\bigcup_{x\in B,B\in \mathcal{T}}\mathcal{D}_{B}(x,i)\right).$$
For $\alpha\in S$, the derived design at $\alpha$ has the block set
$$\mathcal{F}(\alpha)=\bigcup_{B\in \mathcal{B}}\mathcal{A}_{B}(\alpha).$$



(1) For any $x\in G$, $G\in \mathcal{G}$ and $i \in I_m$, consider the coloring of the derived design $\mathcal{F}(x, i)$, which forms the block set of an $\mathrm{STS}(mgn+s-1, mg+s-1)$ on $Y \setminus\{(x, i)\}$ with a hole $((G \times I_m)\cup S) \setminus\{(x,i)\} $.

For any $y\in X\setminus G$ and $1\leq j \leq \frac{m}{2}$, let
$$\mathcal{F}^{y,j}(x,i) = \mathcal{A}_{B_{1}}^{y,j}(x,i) \bigcup \left( \bigcup _{\{x,y\}\subset B,B\in \mathcal{T}}\mathcal{D}_{B}^{y,j}(x,i)\right),$$
where $B_{1}$ is the unique block of $\mathcal{B}$ such that $\{x,y\}\subset B_{1}$.
For any $y\in G\setminus \{x\}$ and $1\leq j \leq \frac{m}{2}$, let
$$\mathcal{F}^{y,j}(x,i) =  \bigcup _{\{x,y\}\subset B,B\in \mathcal{T}}\mathcal{D}_{B}^{y,j}(x,i).$$
For $1\leq j \leq  \frac{m+s}{2}$, let
$$\mathcal{F}^{j}(x,i) = \bigcup_{x\in B, B \in \mathcal{B}}\mathcal{A}_{B}^{j}(x,i).$$
Then it follows that the collection $\{\mathcal{F}^{y,j}(x,i):y\in X\setminus \{x\},1\leq j \leq \frac{m}{2}\}\cup \{\mathcal{F}^{j}(x,i):1\leq j \leq  \frac{s+m}{2}\}$ forms a resolution of $\mathcal{F}(x,i)$, where each $\mathcal{F}^{y,j}(x,i)$ with $y\in X\setminus G,1\leq j \leq \frac{m}{2}$ is a $\mathrm{PPC}$ of $Y\setminus \{(x,i)\}$ and those $\mathcal{F}^{y,j}(x,i)$ with $y\in G\setminus \{x\},1\leq j \leq \frac{m}{2}$ and $\mathcal{F}^{j}(x,i)$ with $1\leq j \leq  \frac{m+s}{2}$ are $\mathrm{PPC}$s of $Y\setminus((G \times I_m) \cup S)$.



(2) Consider the coloring of the derived design $\mathcal{F}(\alpha)$, $\alpha\in S$, which forms the block set of a $\mathrm{GDD}(2,3,mgn)$ of type $(mg)^{n}$ on $X\times I_{m}$ with group set $\mathcal{H}$.
It is obvious that it is a $\mathrm{KF}((mg)^{n})$ with a resolution $\{\mathcal{F}^{y,j}(\alpha):y\in X,1\leq j\leq \frac{m}{2}\}$, where
$$\mathcal{F}^{y,j}(\alpha) =  \bigcup_{y\in B,B\in \mathcal{B}}\mathcal{A}_{B}^{y,j}(\alpha).$$
Hence, we end up with a $\mathrm{gcDCQS}((mg)^n: s)$.
\qed

{\noindent\textbf{Remark:}} From the proof of Construction \ref{C2}, for the case of $mgn+s\equiv 2\pmod 6$, one can weaken the restrictions to $mgn+s\equiv mk_{1}+s\equiv 2\pmod 6$ for some $k_{1}\in K_{1}$.


The rest of this subsection introduces a special combinatorial structure denoted $2$-$\mathrm{RSQS}^{*}(v)$, which is proposed in \cite{DRSQS11}. This design will be used to present another recursive construction for $\mathrm{gcDCQS}$s.

Let $X$ be a set of $v$ elements. A collection $\{(X,\mathcal{B}_{k}^{l}):1\leq k \leq \frac{v-1}{3},1\leq l \leq 3\}$ of $v-1$ $\mathrm{S}(2,4,v)$s is denoted by $2$-$\mathrm{RSQS}^{*}(v)$ if this collection has the following properties:\vspace{-0.2cm}
\begin{itemize}
  \item [(i)] for $1\leq k \leq \frac{v-1}{3}$, the three block sets $\mathcal{B}_{k}^{l}$ $(1\leq l \leq 3)$ share a common $\mathrm{PC}$ of $X$, say $P_{k}$;\vspace{-0.2cm}
  \item [(ii)] $(X,\mathcal{B}^{'})$, where $\mathcal{B}^{'} = \bigcup_{1\leq k \leq \frac{v-1}{3}}P_{k}$, is an $\mathrm{S}(2,4,v)$ (called the special $\mathrm{RS}(2,4,v)$);\vspace{-0.2cm}
  \item [(iii)] $(X,\mathcal{B})$, where $\mathcal{B}$ is the usual union of $\mathcal{B}_{k}^{l}$ ($1\leq k \leq \frac{v-1}{3}$, $1\leq l \leq 3$) is an $\mathrm{SQS}(v)$ (called the underlying $\mathrm{SQS}(v)$);\vspace{-0.2cm}
  \item [(iv)] each quadruple of $\mathcal{B}\setminus \mathcal{B}^{'}$ occurs twice in the multiset $\bigcup_{1\leq k \leq \frac{v-1}{3},1\leq l \leq 3}\mathcal{B}_{k}^{l}$, while each quadruple of $\mathcal{B}^{'}$ occurs three times.
\end{itemize}


{\Construction \label{C4} Let $4m+s\equiv 2,4\pmod 6$. Suppose that there exists a $2$-$\mathrm{RSQS}^{*}(v)$. If there exist a $\mathrm{gcDCQS}(m^{4}:s)$ and an $\mathrm{RDTD}(3,4,m)$, then there exists a $\mathrm{gcDCQS}(m^{v}:s)$.}
\proof Note that $v\equiv 4\pmod {12}$ is the necessary condition for the existence of a $2$-$\mathrm{RSQS}^{*}(v)$. Then one readily checks that $mv+s\equiv4m+s\equiv 2,4\pmod 6$. Let the collection $\{(X,\mathcal{B}_{k}^{l}):1\leq k \leq \frac{v-1}{3},1\leq l \leq 3\}$ be the given $2$-$\mathrm{RSQS}^{*}(v)$.
Let $(X,\mathcal{B})$ be the underlying $\mathrm{SQS}(v)$ and $\mathcal{B}^{'}$ be the spacial $\mathrm{RS}(2,4,v)$.
Now, we consider the derived designs of this $2$-$\mathrm{RSQS}^{*}(v)$. For any $x\in X$,  we denote the derived design at $x$ of $\mathcal{B}$, $\mathcal{B}_{k}^{l}$ and $P_{k}$ by $\mathcal{B}_{x}$, $\mathcal{B}_{x,k}^{l}$ and $P_{x,k}$, respectively. They have the following properties:\vspace{-0.2cm}
\begin{itemize}
  \item [$\bullet$] for $1\leq k \leq \frac{v-1}{3}$, the three block sets $\mathcal{B}_{x,k}^{l}$ $(1\leq l \leq 3)$ share a common block $P_{x,k}$;\vspace{-0.2cm}
  \item [$\bullet$] $\mathcal{B}_{x}^{'}$, where $\mathcal{B}_{x}^{'} = \bigcup_{1\leq k \leq \frac{v-1}{3}}P_{x,k}$, is a $\mathrm{PC}$ of $X\setminus\{x\}$;\vspace{-0.2cm}
  \item [$\bullet$] each triple of $\mathcal{B}_{x}\setminus \mathcal{B}_{x}^{'}$ occurs twice in the multiset $\bigcup_{1\leq k \leq \frac{v-1}{3},1\leq l \leq 3}\mathcal{B}_{x,k}^{l}$, while each triple of $\mathcal{B}_{x}^{'}$ occurs three times.\vspace{-0.2cm}
\end{itemize}

Let $S=\{\alpha_{1},\alpha_{2},\ldots,\alpha_{s}\}$, $Y = (X\times I_{m})\cup S$ and $\mathcal{G}_{B} =\{\{x\}\times I_{m}:x\in B\}$, where $B\subseteq X$. We will construct the desired $\mathrm{gcDCQS}(m^{v}:s)$ on the point set $Y$ with stem $S$ and group set $\mathcal{G}_{X}$.

For each block $B\in \mathcal{B}\setminus\mathcal{B}^{'}$, we construct an $\mathrm{RDTD}(3,4,m)$ $(B\times I_{m}, \mathcal{G}_{B},\mathcal{A}_{B})$. For every point $(x,i)\in B\times I_{m}$, denote the corresponding derived design as $((B\setminus\{x\})\times I_{m}, \mathcal{G}_{B\setminus\{x\}},\mathcal{A}_{B}(x,i))$ and each $\mathcal{A}_{B}(x,i)$ has a resolution $\{\mathcal{A}_{B}^{j}(x,i):1\leq j \leq m\}$, where each $\mathcal{A}_{B}^{j}(x,i)$ is a $\mathrm{PC}$ of $(B\setminus \{x\})\times I_{m}$.

For each block $B\in\mathcal{B}^{'}$, we construct a $\mathrm{gcDCQS}(m^{4}:s)$ $((B\times I_{m})\cup S, S, \mathcal{G}_{B},\mathcal{C}_{B})$. Let
\begin{align*}
\delta = \begin{cases}
\frac{m+s}{2}, & \text{ if } 4m+s\equiv  2 \pmod{6},\\
\frac{m+s}{2}-1, & \text{ if } 4m+s\equiv 4 \pmod{6}.\\
		   \end{cases}
\end{align*}
For every point $(x,i)\in B\times I_{m}$, denote the corresponding derived design as $(((B\times I_{m})\cup S)\setminus\{(x,i)\}, ((\{x\}\times I_{m})\cup S)\setminus\{x,i\},\mathcal{C}_{B}(x,i))$ and each $\mathcal{C}_{B}(x,i)$ has a resolution $\{\mathcal{C}_{B}^{j}(x,i):1\leq j \leq \frac{3m}{2}+\delta\}$, where each $\mathcal{C}_{B}^{j}(x,i)$ is a $\mathrm{PPC}$ of $((B\times I_{m})\cup S)\setminus\{(x,i)\}$; particularly, those $\mathcal{C}_{B}^{j}(x,i)$ with $\frac{3m}{2} + 1 \leq j \leq \frac{3m}{2}+\delta$ are $\mathrm{PPC}$s of $(B\setminus \{x\})\times I_{m}$.
For every point $\alpha\in S$, denote the corresponding derived design as $(B\times I_{m},\mathcal{G}_{B},\mathcal{C}_{B}(\alpha))$ and each $\mathcal{C}_{B}(\alpha)$ has a resolution $\{\mathcal{C}_{B}^{\eta,j}(\alpha):1\leq j \leq \frac{m}{2},\eta\in B\}$, where each $\mathcal{C}_{B}^{\eta,j}(\alpha)$ is a $\mathrm{PC}$ of $(B\setminus\{\eta\})\times I_{m}$.

Let
$$\mathcal{F} = \bigg(\bigcup_{B\in \mathcal{B}\setminus\mathcal{B}^{'}}\mathcal{A}_{B}\bigg)\bigcup \bigg(\bigcup_{B\in \mathcal{B}^{'}}\mathcal{C}_{B}\bigg).$$%\bigcup \bigg(\bigcup_{x\in X}\mathcal{D}_{x}\bigg).$$
Clearly, $(Y,S,\mathcal{G}_{X},\mathcal{F})$ is the block set of a $\mathrm{CQS}(m^{v}:s)$. For $(x,i)\in X\times I_{m}$, the derived design at $(x,i)$ has the block set
$$\mathcal{F}_{(x,i)} = \bigg(\bigcup_{x\in B,B\in \mathcal{B}\setminus\mathcal{B}^{'}}\mathcal{A}_{B}(x,i)\bigg)\bigcup \bigg(\bigcup_{x\in B,B\in \mathcal{B}^{'}}\mathcal{C}_{B}(x,i)\bigg).$$%\bigcup\mathcal{D}_{x}(x,i).$$
Next, we check that each block set $\mathcal{F}_{(x,i)}$ can be partitioned into $\frac{m(v-1)}{2}$ $\mathrm{PPC}$s of $Y\setminus\{(x,i)\}$ and $\delta$ $\mathrm{PPC}$s of $(X\setminus\{x\})\times I_{m}$.

For $1\leq k \leq \frac{v-1}{3}$, $1\leq l \leq 3$ and $1\leq r \leq \frac{m}{2}$, let
$$\mathcal{F}_{(x,i)}(k,l,r) = \bigg(\bigcup_{x\in B,B\setminus\{x\}\in \mathcal{B}_{x,k}^{l}\setminus \{P_{x,k}\}}\mathcal{A}_{B}^{r^{\prime}}(x,i)\bigg)\bigcup \mathcal{C}_{P_{x,k}\cup\{x\}}^{\frac{m}{2}(l-1)+r}(x,i),$$
where the value $r^{\prime}$ equals $r$ or $r+\frac{m}{2}$ according to the first or the second occurrence of $B$ in $\mathcal{B}\setminus\mathcal{B}^{'}$.

For $1 \leq j \leq \delta$, let
$$\mathcal{F}_{(x,i)}(j) = \bigcup_{P_{x,k}\in \mathcal{B}_{x}^{'},1\leq k \leq \frac{v-1}{3}}\mathcal{C}_{P_{x,k}\cup\{x\}}^{\frac{3m}{2}+j}(x,i).$$%\bigcup\mathcal{D}_{x}^{j}(x,i).$$
So each $\mathcal{F}_{(x,i)}$ has a  desired resolution $\{\mathcal{F}_{(x,i)}(k,l,r):1\leq k \leq \frac{v-1}{3},1\leq l \leq 3,1\leq r \leq \frac{m}{2}\}\bigcup\{\mathcal{F}_{(x,i)}(j):1 \leq j \leq \delta\}$.

For $\alpha\in S$, the derived design at $\alpha$ has the block set
$$\mathcal{F}_{\alpha}=\bigcup_{B\in \mathcal{B}^{\prime}}\mathcal{C}_{B}(\alpha).$$
%It is left to check that each block set $\mathcal{F}_{\alpha}$ can be partitioned into $\frac{mv}{2}$ $\mathrm{PC}$s of $(X\setminus \eta)\times I_{m}$ for some $\eta\in X$.
For $\eta\in X$, $1 \leq j \leq \frac{m}{2}$, let
$$\mathcal{F}_{\alpha}(\eta,j)=\bigcup_{\eta\in B,B\in \mathcal{B}^{\prime}}\mathcal{C}_{B}^{\eta,j}(\alpha).$$
So each $\mathcal{F}_{\alpha}$ is a block set of a $\mathrm{KF}(m^{v})$.
Then we obtain a $\mathrm{gcDCQS}(m^{v}:s)$.
\qed

\begin{Theorem}\label{2RDSQS}
\begin{itemize}
  \item [$(1)$] {\cite[Corollary 3.2]{DRSQS11}} There is a $2$-$\mathrm{RSQS}^{*}(4^{n})$ for any integer $n\geq 2$.
 \item [$(2)$] {\cite{J_TD,DRSQS11}} There is an $\mathrm{RDTD}(3,4,m)$ for any positive integer $m$ satisfying $\gcd(m,4) \neq 2$ and $\gcd(m,18) \neq 3$.
 \end{itemize}
\end{Theorem}



\begin{Corollary}\label{Rmul}
Let $m\equiv 4,8\pmod {12}$. If there exist a $\mathrm{gcDCQS}(m^{4}:0)$ and an $\mathrm{mcDSQS}(m)$, then there exists an $\mathrm{mcDSQS}(m\cdot 4^{n})$ for any integer $n\geq 2$.
\end{Corollary}
\proof Let $m\equiv 4,8\pmod {12}$ and $s=0$. For any integer $n\geq 2$, apply Construction \ref{C4}. Since a $2$-$\mathrm{RSQS}^{*}(4^{n})$ and an $\mathrm{RDTD}(3,4,m)$ exist by Theorem \ref{2RDSQS}, if there exists a $\mathrm{gcDCQS}(m^{4}:0)$, then there exists a $\mathrm{gcDCQS}(m^{4^{n}}:0)$. Further, if there exists an $\mathrm{mcDSQS}(m)$,
then there exists an $\mathrm{mcDSQS}(m\cdot 4^{n})$ by Construction \ref{mcdcqs-mcdsqs}.
\qed

\noindent{\bf Remark:} Construction \ref{C4} presents a new construction even for $\mathrm{RDCQS}$s. As a special case, taking $m=4$ and $s=0$ in Construction \ref{C4} gives the quadrupling construction by Lui and Lei \cite[Construction 3.2]{LEI1} (noting the existence of an RDCQS$(4^4:0)$ and an $\mathrm{RDTD}(3,4,4)$).



\subsection{Special considerations}

Construction \ref{mcdcqs-mcdsqs} provides a construction of an $\mathrm{mcDSQS}(gn+s)$ by filling in holes of a $\mathrm{gcDCQS}(g^{n}:s)$. Constructions \ref{C3}-\ref{C4} present approaches of producing $\mathrm{gcDCQS}$s. However, in order to deal with $\mathrm{mcDSQS}(v)$ where $v\equiv 2\pmod 6$, we have to construct some input designs, say $\mathrm{gcDCQS}(g^{k}:s)$ with $g+s\geq 20$ (because $\chi^{\prime}(v) >\frac{v+1}{2}$ if $v\in\{7,13\}$), which requires strong computing power.
In this subsection, we consider $\mathrm{CQS}(g^{n}:s)$ with group size $g \equiv2\pmod 6$ and stem size $s=2$. We introduce two new types of colorings of the derived designs and define $\mathrm{^{1}cDCQS}(g^{n}:2)$ and $\mathrm{^{2}cDCQS}(g^{n}:2)$; $\mathrm{^{1}cDCQS}$ is used in a new filling construction %with $\mathrm{RDSQS}(g+2)$ to produce $\mathrm{mcDSQS}(gn+2)$
and $\mathrm{^{2}cDCQS}$ is an instrumental design to generate $\mathrm{^{1}cDCQS}$s in a recursive construction.


Let $(X,\mathcal{G},\mathcal{B})$ be a $\mathrm{GDD}(2,3,gn)$ of type $g^{n}$ with $g$ even. We say that $(X,\mathcal{G},\mathcal{B})$ is {\em{good colorable}} and denoted by $\mathrm{gcGDD}(2,3,gn)$ of type $g^{n}$ if the block set $\mathcal{B}$ can be partitioned into $\frac{gn}{2}+1$ $\mathrm{PPC}$s of $X\setminus G$, $G\in \mathcal{G}$, to be specific,  $\mathcal{B} = \left(\bigcup_{G\in \mathcal{G},1\leq i \leq \frac{g}{2}} P_{G}^{i}\right)\bigcup P_{H}^{*}$, where $P_{G}^{i}$ is a $\mathrm{PPC}$ of $X \setminus G$ and $P_{H}^{*}$ is a $\mathrm{PPC}$ of $X\setminus H$ for some $H\in \mathcal{G}$. We call $H$ the special group of the $\mathrm{gcGDD}(2,3,gn)$ of type $g^{n}$.

Now, we define two coloring types for the derived designs of a $\mathrm{CQS}(g^{n}:2)$. Let $(X, S, \mathcal{G}, \mathcal{A})$ be a $\mathrm{CQS}(g^n: 2)$. For $x \in X$, put $\mathcal{A}_x=\{A \setminus\{x\}: x \in A, A \in \mathcal{A}\}$.
A $\mathrm{CQS}(g^n: 2)$ whose derived designs admit type I colorings, denoted by $^{1}\mathrm{cDCQS}(g^n: 2)$, refers to a $\mathrm{CQS}(g^n: 2)$ whose derived designs $\mathcal{A}_x$, $x \in X$, have the following properties:\vspace{-0.2cm}
\begin{itemize}
  \item [(1)] for $x\in G$ and $G\in \mathcal{G}$, each $(X\setminus \{x\}, (G\cup S)\setminus \{x\},\mathcal{A}_{x})$ is a $\mathrm{gcSTS}(gn+1,g+1)$;\vspace{-0.2cm}
  \item [(2)] for $x\in S$, each $(X\setminus S, \mathcal{G}, \mathcal{A}_{x})$ is a $\mathrm{gcGDD}(2,3,gn)$ of type $g^{n}$ with the same special group.\vspace{-0.2cm}
\end{itemize}
A $\mathrm{CQS}(g^n: 2)$ whose derived designs admit type II colorings, denoted by $^{2}\mathrm{cDCQS}(g^n: 2)$, refers to a $\mathrm{CQS}(g^n: 2)$ whose derived designs $\mathcal{A}_x$, $x \in X$, have the following properties:\vspace{-0.2cm}
\begin{itemize}
  \item [(1)] for $x\in G$ and $G\in \mathcal{G}$, each $(X\setminus \{x\}, (G\cup S)\setminus \{x\},\mathcal{A}_{x})$ is a $\mathrm{gcSTS}(gn+1,g+1)$;\vspace{-0.2cm}

  \item [(2)] for $x\in S$, each $(X\setminus S, \mathcal{G}, \mathcal{A}_{x})$ is an $\mathrm{RGDD}(2,3,gn)$ of type $g^{n}$.
\end{itemize}




{\Construction \label{cqs-wrdsqs1} Let $n\equiv 0\pmod 3$. If there exist a $\mathrm{^{1}cDCQS}(g^n: 2)$ and an $\mathrm{RDSQS}(g+2)$, then there exists an $\mathrm{mcDSQS}(gn+2)$.}
\proof Suppose $(X,S,\mathcal{G},\mathcal{A})$ is the given $^{1}\mathrm{cDCQS}(g^n: 2)$.
For $x \in X$, let $\mathcal{A}_x=\{A \setminus\{x\}: x \in A, A \in \mathcal{A}\}$.
Then, each $(X\setminus \{x\}, (G\cup S)\setminus \{x\},\mathcal{A}_{x})$ with $x\in G$ and $G\in \mathcal{G}$ is a $\mathrm{gcSTS}(gn+1,g+1)$ with a hole $(G\cup S)\setminus\{x\}$; its resolution is $\{P_{x}^{j}:1 \leq j \leq \frac{gn}{2}+1\}$, where each $P_{x}^{j}$ is a $\mathrm{PPC}$ of $X\setminus \{x\}$; particularly, those $P_{x}^{j}$ with $1 \leq j \leq \frac{g}{2}+1$ are $\mathrm{PPC}$s of $X\setminus (G\cup S)$.
In addition, each $(X\setminus S, \mathcal{G}, \mathcal{A}_{x})$ with $x\in S$ is a $\mathrm{gcGDD}(2,3,gn)$ of type $g^{n}$ with the special group $H\in \mathcal{G}$; its resolution is $\{P_{x}^{j}(G):G\in \mathcal{G},1\leq j \leq \frac{g}{2}\}\cup\{P_{x}^{*}\}$, where each $P_{x}^{j}(G)$, $G\in \mathcal{G}$ and $1\leq j \leq \frac{g}{2}$, is a $\mathrm{PPC}$ of $X\setminus (G\cup S)$ and $P_{x}^{*}$ is a $\mathrm{PPC}$ of $X\setminus (H\cup S)$.

For each $G\in \mathcal{G}$, we can construct by assumption an $\mathrm{RDSQS}(g+2)$ on $G\cup S$ with block set $\mathcal{B}(G)$.
So we have $g+2$ derived $\mathrm{KTS}(g+1)$s, say $((G\cup S) \setminus\{x\}, \mathcal{B}_{x}(G))$ with $x\in G\cup S$.
Further, each $\mathcal{B}_{x}(G)$ with $x\in G\cup S$ has a resolution $\{Q_{x}^{j}(G):1\leq j \leq \frac {g}{2}\}$ into $\mathrm{PC}$s of $(G\cup S)\setminus\{x\}$.
Define
$$\mathcal{F} = \mathcal{A}\bigcup\left(\bigcup_{G\in\mathcal{G}}\mathcal{B}(G)\right).$$
Then, $\mathcal{F}$ is the block set of an $\mathrm{SQS}(gn+2)$.
For $x\in G$ and $G\in \mathcal{G}$, the derived design at $x$ has the block set $\mathcal{A}_{x}\cup \mathcal{B}_{x}(G)$ and a resolution
$$\big\{P_{x}^{j}\cup Q_{x}^{j}(G):1\leq j \leq \frac{g}{2}\big\}\bigcup\big\{P_{x}^{j}:\frac{g}{2}+1 \leq j \leq  \frac{gn}{2}+1\big\}.$$
For $x\in S$, the derived design at $x$ has the block set $\mathcal{A}_{x}\bigcup (\bigcup_{G\in \mathcal{G}}\mathcal{B}_{x}(G))$ and a resolution
$$\big\{P_{x}^{j}(G)\cup Q_{x}^{j}(G):G\in \mathcal{G},1\leq j \leq \frac{g}{2}\big\}\bigcup\big\{P_{x}^{*}\big\}.$$
Hence, we end up with the desired $\mathrm{mcDSQS}(gn+2)$.
\qed





{\Construction \label{C1} Let $m\equiv 2 \pmod 6$. Suppose that there exists an $\mathrm{RDS}(3,k+1,n+1)$. If there exist a $\mathrm{^{1}cDCQS}(m^{k}: 2)$, a $\mathrm{^{2}cDCQS}(m^{k}: 2)$ and an $\mathrm{RDGDD}(3,4, (k+1)\{m\})$, then there exists a $\mathrm{^{1}cDCQS}(m^n: 2)$.}

\proof Suppose $(I_n \cup\{\infty\}, \mathcal{B})$ is the given $\mathrm{RDS}(3, k+1, n+1)$. Let $\mathcal{B}_\infty=\{B \backslash\{\infty\}: \infty\in B, B \in \mathcal{B}\}$ and $\mathcal{T}=\{B: \infty \notin B, B \in \mathcal{B}\}$. Thus $(I_n, \mathcal{B}_{\infty},\mathcal{T})$ forms a $1$-fan $\mathrm{S}(3,(\{k\}, \{k+1\}), n)$, where $(I_{n},\mathcal{B}_{\infty})$ is an $\mathrm{S}(2,k,n)$. Let $S=\{\alpha_1, \alpha_2\}, Y=(I_n \times I_m) \cup S$, $\mathcal{H}=\{\{x\} \times I_m: x \in I_n\}$ and $\mathcal{H}_{B} = \{\{x\}\times I_{m}:x\in B\}$, where $B\subseteq I_{n}$.

Fix an element $\eta\in I_{n}$. For any $B\in \mathcal{B}_{\infty}$ and $\eta \in B$, we can construct by assumption a $^{1}\mathrm{cDCQS}(m^{k}: 2)$ $((B\times I_{m})\cup S,S,\mathcal{H}_{B},\mathcal{A}_B)$ with derived designs $\mathcal{A}_{B}(x,i)$ with $(x,i)\in B\times I_{m}$ and $\mathcal{A}_{B}(\alpha)$ with $\alpha\in S$. Moreover, the following should be satisfied:\vspace{-0.2cm}
\begin{itemize}
  \item [(A1)] each $\mathcal{A}_{B}(x, i)$, where $(x, i) \in B \times I_m$, is the block set of a $\mathrm{gcSTS}(mk+1, m+1)$ on $((B \times I_m) \setminus\{(x, i)\}) \cup S$ with a hole $((\{x\} \times I_m)\cup S) \setminus\{(x,i)\}$;
       its resolution is $\{\mathcal{A}_{B}^{j}(x, i):1\leq j \leq \frac{mk}{2}+1\}$, where each $\mathcal{A}_{B}^{j}(x, i)$ is a $\mathrm{PPC}$ of $((B \times I_m) \setminus\{(x, i)\}) \cup S$. Particularly,
       those $\mathcal{A}_{B}^{j}(x, i)$ with $\frac{m(k-1)}{2} +1 \leq j\leq \frac{mk}{2}+1$ are $\mathrm{PPC}$s of  $(B \setminus\{x\})\times I_{m}$;\vspace{-0.2cm}

  \item [(A2)] each $(B\times I_{m},\mathcal{H}_{B},\mathcal{A}_B(\alpha))$, where $\alpha\in S$, is a $\mathrm{gcGDD}(2,3,mk)$ of type $m^{k}$ with the special group $\{\eta\}\times I_{m}$; and $\mathcal{A}_B(\alpha)$ has a resolution $\{\mathcal{A}_{B}^{y,j}(\alpha):y\in B, 1\leq j\leq \frac{m}{2}\}\cup\{\mathcal{A}_{B}^{\eta,0}(\alpha)\}$, where each $\mathcal{A}_{B}^{y,j}(\alpha)$ is a $\mathrm{PPC}$ of $(B \setminus\{y\})\times I_{m}$;\vspace{-0.2cm}
\end{itemize}

For any $B\in \mathcal{B}_{\infty}$ and $\eta \notin B$, we can construct by assumption a $\mathrm{^{2}cDCQS}(m^{k}: 2)$ $((B\times I_{m})\cup S,S,\mathcal{H}_{B},\mathcal{A}_B)$; its derived designs are $\mathcal{A}_{B}(x,i)$ with $(x,i)\in B\times I_{m}$ and $\mathcal{A}_{B}(\alpha)$ with $\alpha\in S$. Moreover, each $\mathcal{A}_{B}(x, i)$, $(x, i) \in B \times I_m$, satisfies the property (A$1$). In addition, each $(B\times I_{m},\mathcal{H}_{B},\mathcal{A}_B(\alpha))$ with $\alpha\in S$ is an $\mathrm{RGDD}(2,3,mk)$ of type $m^{k}$ with a resolution $\{\mathcal{A}_{B}^{j}(\alpha):1\leq j \leq \frac{m(k-1)}{2}\}$ into $\mathrm{PC}$s of $B\times I_{m}$.


For any $B\in \mathcal{T}$, we can construct by assumption an $\mathrm{RDGDD}(3,4,(k+1)\{m\})$ $(B\times I_{m},\mathcal{H}_{B},\mathcal{D}_B)$; its derived designs are $\mathcal{D}_{B}(x,i)$ with $(x,i)\in B\times I_{m}$.
Moreover, each $\mathcal{D}_B(x, i)$, where $(x, i) \in B \times I_m$ and $B \in \mathcal{T}$, forms the block set of an $\mathrm{RGDD}(2,3, mk)$ with group set $\mathcal{H}_{B\setminus\{x\}}$; its resolution is $\{\mathcal{D}_B^j(x, i): 1 \leqslant j \leqslant \frac{m(k-1)}{2}\}$, each $\mathcal{D}_B^j(x, i)$ being a $\mathrm{PC}$ of $(B \setminus\{x\}) \times I_m$.
Define
$$\mathcal{F} = \left(\bigcup_{B\in \mathcal{B}_{\infty}}\mathcal{A}_{B}\right)\bigcup\left(\bigcup_{B\in \mathcal{T}}\mathcal{D}_{B}\right).$$
Then $(Y,S,\mathcal{H},\mathcal{F})$ forms a $\mathrm{CQS}(m^{n}:2)$. For $(x,i)\in I_{n}\times I_{m}$, the derived design at $(x,i)$ has the block set
$$\mathcal{F}(x,i) = \left(\bigcup_{x\in B,B\in \mathcal{B}_{\infty}}\mathcal{A}_{B}(x,i)\right)\bigcup\left(\bigcup_{x\in B,B\in \mathcal{T}}\mathcal{D}_{B}(x,i)\right).$$
For $\alpha\in S$, the derived design at $\alpha$ has the block set
$$\mathcal{F}(\alpha)=\bigcup_{B\in \mathcal{B}_{\infty}}\mathcal{A}_{B}(\alpha).$$
By the proof of Construction \ref{C3}, each derived design $\mathcal{F}(x, i)$, $(x, i) \in I_n \times I_m$, forms the block set of a $\mathrm{gcSTS}(mn+1, m+1)$ on $Y \setminus\{(x, i)\}$ with a hole $(\{x\} \times(I_m \setminus\{i\})) \cup S$.



Next, consider the coloring of the derived design $\mathcal{F}(\alpha)$, $\alpha\in S$.
Since $\mathcal{B}_{\infty}$ is a resolvable $\mathrm{S}(2,k,n)$, it has a resolution $\{\mathcal{B}_{\infty}(h): 1 \leqslant h \leqslant p\}$, where $p=(n-1) /(k-1)$. So there is a unique block $C_{\infty}(h)$ in $\mathcal{B}_{\infty}(h)$, $1 \leqslant h \leqslant p$, that contains $\eta$.
For any block $C_{\infty}(h)$, $1 \leqslant h \leqslant p$, we define a bijection $f_{B}$ from $(B\setminus\{\eta\})\times \{1,2,\ldots, \frac{m}{2}\}$ to $\{1,2,\ldots, \frac{m(k-1)}{2}\}$ arbitrarily.
The collection $\{C_{\infty}(h)\setminus \{\eta\}:1 \leqslant h \leqslant p\}$ is a partition of $I_{n}\setminus \{\eta\}$.

For every $y\in I_{n}$, $y\neq\eta$ and $1\leq j\leq \frac{m}{2}$, there is a unique $h$ such that $y\in C_{\infty}(h)$. Let
$$\mathcal{F}^{y,j}(\alpha) = \mathcal{A}_{C_{\infty}(h)}^{y,j}(\alpha)\bigcup \left(\bigcup_{B\in \mathcal{B}_{\infty}(h)\setminus \{C_{\infty}(h)\}}\mathcal{A}_{B}^{f_{C_{\infty}(h)}(y,j)}(\alpha)\right).$$
For $0\leq j\leq \frac{m}{2}$, let
$$\mathcal{F}^{\eta,j}(\alpha) = \bigcup_{B\in \mathcal{B}_{\infty},\eta\in B}\mathcal{A}_{B}^{\eta,j}(\alpha).$$
We can check that the collection $\{\mathcal{F}^{y,j}(\alpha):y\in I_{n},1\leq j\leq \frac{m}{2}\}\bigcup\{\mathcal{F}^{\eta,0}(\alpha)\}$ forms a resolution of $\mathcal{F}(\alpha)$.
In particular,
each $\mathcal{F}^{y,j}(\alpha)$ ($y\in I_{n}$, $y\neq \eta$ and $1\leq j\leq \frac{m}{2}$) is a $\mathrm{PPC}$ of $(I_{n}\setminus\{y\})\times I_{m}$ and each $\mathcal{F}^{\eta,j}(\alpha)$ ($0\leq j\leq \frac{m}{2}$) is a $\mathrm{PPC}$ of $(I_{n}\setminus\{\eta\})\times I_{m}$.
Hence, we end up with a $\mathrm{^{1}cDCQS}(m^n: 2)$.
\qed


\section{Existence and applications}

In this section, we produce new infinite families  of  $\mathrm{RDSQS}(6n+4)$s and  $\mathrm{mcDSQS}(6n+2)$s. Then  we show their applications in both
large sets of Kirkman triple systems and diameter perfect constant-weight codes.
%We construct an $\mathrm{RDCQS}(8^{4}:2)$ required in Constructions \ref{mcdcqs-mcdsqs} and \ref{C4} and establish the existence of an infinite family of $\mathrm{RDSQS}(2^{2n+1}+2)$. Further, some new infinite families of $\mathrm{LKTS}$s are given.
%By direct constructions we display the existence of an $\mathrm{mcDSQS}(v)$ for $v=26,32$.
%We construct some necessary ingredient
%designs required in Constructions \ref{cqs-wrdsqs1} and \ref{C1} and establish the existence of an infinite family of $\mathrm{mcDSQS}(2\cdot 9^{n}+2)$.
%As a consequence, we show their applications in diameter perfect constant-weight codes.

{\Lemma \label{CQS} There exists an $\mathrm{RDCQS}(8^{4}:2)$.}
\proof Let $X=\mathbb{Z}_{32}\cup\{a,b\}$ and $\mathcal{G} = \{G_i: 0\leq i \leq 3\}$, where $G_i=\{i, 4+i, 8+i, 12+i, 16+i, 20+i, 24+i, 28+i\}$.
We construct a $\mathrm{CQS}(8^{4}:2)$ on point set $X$ with stem $S=\{a,b\}$ and group set $\mathcal{G}$.
Let $\mathcal{H}$ be the permutation group generated by the  permutation
%\begin{align*}
$(0,1,\ldots,31)(a,b)$.
%\end{align*}
Our $\mathrm{CQS}(8^{4}:2)$ $(X,S,\mathcal{G},\mathcal{B})$ admits the automorphism group $\mathcal{H}$ and the base blocks are listed as follows:
\begin{center}
\begin{tabular}{l l l l l l l l }
\{0, 1, 2, $a$\}& \{0, 1, 3, 16\}& \{0, 1, 4, 13\}& \{0, 1, 5, 6\}& \{0, 1, 7, 9\}\\
\{0, 1, 8, 23\}& \{0, 1, 10, 25\}& \{0, 1, 11, 12\}& \{0, 1, 14, 19\}& \{0, 1, 15, 18\}\\
\{0, 1, 17, 30\}& \{0, 1, 20, 29\}& \{0, 1, 24, 26\}& \{0, 2, 4, 18\}& \{0, 2, 5, 24\}\\
\{0, 2, 6, 28\}& \{0, 2, 7, 23\}& \{0, 2, 10, 29\}& \{0, 2, 11, 27\}& \{0, 2, 12, 21\}\\
\{0, 2, 13, 22\}& \{0, 2, 14, 20\}& \{0, 2, 17, $b$\}& \{0, 3, 6, $a$\}& \{0, 3, 7, 20\}\\
\{0, 3, 8, 25\}& \{0, 3, 9, 26\}& \{0, 3, 10, 27\}& \{0, 3, 11, 14\}& \{0, 3, 15, 28\}\\
\{0, 4, 9, 27\}& \{0, 4, 11, 15\}& \{0, 4, 14, 22\}& \{0, 5, 10, $a$\}& \{0, 5, 11, 26\}\\
\{0, 5, 12, 17\}& \{0, 6, 12, 22\}& \{0, 6, 13, 24\}& \{0, 6, 14, 25\}& \{0, 6, 19, $b$\}\\
\{0, 7, 14, $a$\}& \{0, 9, 18, $b$\}& \{0, 10, 21, $b$\}\\
\end{tabular}
\end{center}
The resolutions of derived designs at $0$ and $a$ are listed in Appendix A. Thus, we obtain an
$\mathrm{RDCQS}(8^{4}:2)$.
\qed

{\Theorem\label{result_2} There exists an $\mathrm{RDSQS}(2^{2n+1}+2)$ for any nonnegative integer $n$.}
\proof For $n=0$, an $\mathrm{RDSQS}(4)$ exists obviously. For $n=1$, it is well-known that an $\mathrm{SQS}(10)$ is unique up to isomorphism and its derived designs are all resolvable.
For $n=2$, apply Construction \ref{mcdcqs-mcdsqs}. Since an $\mathrm{RDCQS}(8^{4}:2)$ exists by Lemma \ref{CQS},  an $\mathrm{RDSQS}(34)$ exists.

For $n\geq 3$, apply Construction \ref{C4} with a $2$-$\mathrm{RSQS}^{*}(4^{n-1})$, which exists by Theorem \ref{2RDSQS}. Since an $\mathrm{RDCQS}(8^{4}:2)$ and an $\mathrm{RDTD}(3,4,8)$ exist by Lemma \ref{CQS} and Theorem \ref{2RDSQS},  an $\mathrm{RDCQS}(8^{4^{n-1}}:2)$ exists. Further, apply Construction \ref{mcdcqs-mcdsqs} to obtain an $\mathrm{RDSQS}(2^{2n+1}+2)$.
\qed


\begin{Corollary}\label{result_4}
There exists  an $\mathrm{LKTS}(3\cdot 2^{2n+1}+3)$ for  any nonnegative integer $n$.
\end{Corollary}
\proof Apply Theorems \ref{RDS_1}  and \ref{result_2}. 	
\qed


%\begin{Lemma}[\label{RDS_3}{\cite{yuan_olkts}}]
%If there exists an $\mathrm{RDSQS}(u + 1)$ and an $\mathrm{RDSQS}(v + 1)$, then there exists
%an $\mathrm{RDSQS}(uv + 1)$.
%\end{Lemma}
%
%\begin{Theorem}[\label{RDS_4}{\cite{LEI1}}]
%Let $L =\{7\cdot4^{n}-1 : n \geq 1\} \cup \{3\cdot7^{n} : n \geq 0\} \cup \{3\cdot13^{n} : n \geq 0\} \cup \{4^{m}25^{n}-1 : m \geq 1, n \geq 0\} \cup \{4 \cdot 7^{n} - 1 : n \geq 0\} \cup \{2q^{n} + 1 : n \geq 0, q \equiv 7 \pmod {12} \text{ and } q \text{ is
%a prime power}\} \cup \{45, 57, 81, 93, 105, 117, 141, 165\}$. Then there exists an $\mathrm{RDSQS}(v + 1)$
%and an $\mathrm{LKTS}(3v)$ where $v$ is the product of some elements in $L$.
%\end{Theorem}
%
%\begin{Theorem}\label{result_4}
%Let $L =\{2^{2n+1}+1:n\geq 0\}\cup\{7\cdot4^{n}-1 : n \geq 1\} \cup \{3\cdot7^{n} : n \geq 0\} \cup \{3\cdot13^{n} : n \geq 0\} \cup \{4^{m}25^{n}-1 : m \geq 1, n \geq 0\} \cup \{4 \cdot 7^{n} - 1 : n \geq 0\} \cup \{2q^{n} + 1 : n \geq 0, q \equiv 7 \pmod {12} \text{ and } q \text{ is
%a prime power}\} \cup \{45, 57, 81, 93, 105, 117, 141, 165\}$. Then there exists an $\mathrm{RDSQS}(v + 1)$
%and an $\mathrm{LKTS}(3v)$ where $v$ is the product of some elements in $L$.
%\end{Theorem}
%\proof Applying Theorems \ref{result_2} and \ref{RDS_4}, to the known product constructions in Lemma \ref{RDS_3}, we obtain an $\mathrm{RDSQS}(v + 1)$s. By Theorem \ref{RDS_1}, there exists an $\mathrm{LKTS}(3v)$. 	
%\qed

{\Lemma\label{MIN_Dsqs} There exists an $\mathrm{mcDSQS}(v)$ for $v=20,26,32$.}
\proof In \cite{SHI}, Shi et al. show that the cyclic $\mathrm{SQS}(20)$ constructed in \cite{CSQS} is an $\mathrm{mcDSQS}(20)$.
We construct a cyclic $\mathrm{mcDSQS}(v)$ on $\mathbb{Z}_{v}$ for $v=26,32$; the base blocks and the resolutions are listed in Appendix B.
\qed

{\Lemma\label{Dcqs20_1} There exists a $\mathrm{^{1}cDCQS}(2^{9}:2)$.}
\proof Let $X=I_{18}\cup\{a,b\}$ and $\mathcal{G} = \{\{i,i+9\}: 0\leq i \leq 8\}$.
We construct a $\mathrm{CQS}(2^{9}:2)$ on the point set $X$ with stem $S=\{a,b\}$ and group set $\mathcal{G}$.
Let $\mathcal{H}$ be a permutation group generated by the following permutation:
$$ (0,3,6,9,12,15)(1,4,7,10,13,16)(2,5,8,11,14,17)(a,b).$$
Our $\mathrm{CQS}(2^{9}:2)$ $(X,S,\mathcal{G},\mathcal{B})$ admits the automorphism group $\mathcal{H}$ and the base blocks are listed as follows. (The underlined blocks belong to short orbits of length $3$ while the double-underlined blocks lie in short orbits of length $2$.)
\begin{center}
\begin{tabular}{l l l l l l l l }
\{0, 1, 2, $a$\}& \{0, 1, 3, 6\}& \{0, 1, 4, 15\}& \{0, 1, 5, 13\}& \{0, 1, 7, 9\} \\
\{0, 1, 8, $b$\}& \{0, 1, 10, 14\}& \{0, 1, 11, 16\}& \{0, 1, 12, 17\}& \{0, 2, 3, 12\}\\
\{0, 2, 4, 5\}& \{0, 2, 6, 15\}& \{0, 2, 7, 17\}& \{0, 2, 8, 16\}& \{0, 2, 10, $b$\}\\
\{0, 2, 13, 14\}& \{0, 3, 8, 14\}& \{0, 3, 10, 13\}& \{0, 3, 11, $b$\}&\{0, 3, 17, $a$\}\\
\{0, 4, 6, 10\}& \{0, 4, 8, 11\}& \{0, 4, 14, $a$\}&\{0, 4, 17, $b$\}& \{0, 5, 8, 12\}\\
\{0, 5, 10, $a$\}& \{0, 5, 14, 17\}& \{0, 5, 16, $b$\}&\{0, 6, 13, $b$\}& \{0, 6, 16, 17\}\\
\{0, 7, 8, 13\}& \{0, 7, 11, $a$\}&\{0, 7, 14, 16\}& \{0, 8, 10, 17\}& \{0, 11, 13, 17\}\\
\{0, 13, 16, $a$\}&\{1, 2, 4, $b$\}&\{1, 2, 5, 14\}& \{1, 2, 13, 16\}&\{1, 4, 11, 17\}\\
\{1, 5, 7, $b$\}& \{1, 7, 14, 17\}&\{2, 5, 8, $b$\}&\underline{\{0, 2, 9, 11\}}& \underline{\{0, 4, 9, 13\}}\\
\underline{\{1, 2, 10, 11\}}& \underline{\{1, 4, 10, 13\}}&\underline{\underline{\{0, 6, 12, $a$\}}}& \underline{\underline{\{1, 7, 13, $a$\}}}& \underline{\underline{\{2, 8, 14, $a$\}}}%& \{2, 11, $a$, $b$\}& \{1, 10, $a$, $b$\}& \{0, 9, $a$, $b$\}
\end{tabular}
\end{center}% 2�����
Then we only need to list the resolution of the derived designs at $0,1,2,a$ and $b$ to check that this $\mathrm{CQS}(2^{9}:2)$ is a $\mathrm{^{1}cDCQS}(2^{9}:2)$. We list below the block set of the derived designs at the point $x\in\{0,1,2\}$; the blocks in each row form a $\mathrm{PPC}$ of $X\setminus\{x\}$ and the last two rows form $\mathrm{PPC}$s of $I_{18}\setminus G$, where $x\in G$, $G\in \mathcal{G}$.
{\small\begin{center}
\begin{tabular}{l l l l l l l l }
$\mathcal{B}_{0}$:
    & \{1, 4, 15\}& \{2, 9, 11\}& \{3, 17, $a$\}& \{5, 8, 12\}& \{6, 13, $b$\}& \{7, 14, 16\}\\
    & \{1, 5, 13\}& \{2, 10, $b$\}& \{3, 4, 7\}& \{6, 16, 17\}& \{8, 15, $a$\}& \{9, 12, 14\}\\
    & \{1, 8, $b$\}& \{2, 3, 12\}& \{4, 14, $a$\}& \{5, 6, 7\}& \{9, 10, 16\}& \{11, 13, 17\}\\
    & \{2, 8, 16\}& \{3, 11, $b$\}& \{4, 9, 13\}& \{5, 14, 17\}& \{6, 12, $a$\}& \{7, 10, 15\}\\
    & \{1, 10, 14\}& \{2, 6, 15\}& \{3, 5, 9\}& \{4, 8, 11\}& \{7, 12, $b$\}& \{13, 16, $a$\}\\
    & \{1, 12, 17\}& \{2, 4, 5\}& \{3, 10, 13\}& \{6, 8, 9\}& \{7, 11, $a$\}& \{14, 15, $b$\}\\
    & \{1, 7, 9\}& \{3, 15, 16\}& \{4, 17, $b$\}& \{5, 10, $a$\}& \{6, 11, 14\}\\
    & \{1, 2, $a$\}& \{5, 16, $b$\}& \{7, 8, 13\}& \{9, 15, 17\}& \{10, 11, 12\}\\
    & \{1, 11, 16\}& \{2, 7, 17\}& \{3, 8, 14\}& \{4, 6, 10\}& \{12, 13, 15\}\\%& \{9, $a$, $b$\}
    & \{1, 3, 6\}& \{2, 13, 14\}& \{4, 12, 16\}& \{5, 11, 15\}& \{8, 10, 17\}\\
\end{tabular}
\end{center}

\begin{center}
\begin{tabular}{l l l l l l l l }
$\mathcal{B}_{1}$: & \{0, 2, $a$\}& \{3, 12, 13\}& \{4, 11, 17\}& \{5, 7, $b$\}& \{6, 10, 15\}& \{8, 14, 16\}\\
    & \{0, 3, 6\}& \{2, 7, 12\}& \{4, 10, 13\}& \{5, 8, 15\}& \{9, 14, $b$\}& \{16, 17, $a$\}\\
    & \{0, 4, 15\}& \{2, 10, 11\}& \{3, 16, $b$\}& \{5, 12, $a$\}& \{6, 13, 14\}& \{8, 9, 17\}\\
    & \{0, 8, $b$\}& \{2, 13, 16\}& \{3, 10, 17\}& \{4, 9, 12\}& \{6, 7, 11\}& \{14, 15, $a$\}\\
    & \{2, 4, $b$\}& \{3, 8, $a$\}& \{5, 6, 17\}& \{7, 10, 16\}& \{9, 13, 15\}& \{11, 12, 14\}\\
    & \{0, 10, 14\}& \{2, 6, 8\}& \{3, 7, 15\}& \{4, 5, 16\}& \{9, 11, $a$\}& \{13, 17, $b$\}\\
    & \{0, 11, 16\}& \{2, 15, 17\}& \{3, 4, 14\}& \{5, 9, 10\}& \{6, 12, $b$\}& \{7, 13, $a$\}\\
    & \{0, 5, 13\}& \{2, 3, 9\}& \{4, 6, $a$\}& \{7, 14, 17\}& \{8, 10, 12\}& \{11, 15, $b$\}\\
    & \{0, 7, 9\}& \{2, 5, 14\}& \{8, 11, 13\}&  \{12, 15, 16\}\\%\{10, $a$, $b$\}&
    & \{0, 12, 17\}& \{3, 5, 11\}& \{4, 7, 8\}& \{6, 9, 16\}\\
\end{tabular}
\end{center}

\begin{center}
\begin{tabular}{l l l l l l l l }
$\mathcal{B}_{2}$:
& \{0, 3, 12\}& \{1, 4, $b$\}& \{5, 6, 11\}& \{8, 14, $a$\}& \{9, 13, 17\}& \{10, 15, 16\}\\
    & \{0, 4, 5\}& \{3, 7, $a$\}& \{6, 12, 17\}& \{8, 10, 13\}& \{9, 16, $b$\}& \{11, 14, 15\}\\
    & \{0, 7, 17\}& \{1, 10, 11\}& \{3, 14, 16\}& \{4, 6, 13\}& \{5, 8, $b$\}& \{9, 12, $a$\}\\
    & \{0, 10, $b$\}& \{1, 5, 14\}& \{3, 8, 17\}& \{4, 11, 12\}& \{6, 7, 16\}& \{13, 15, $a$\}\\
    & \{0, 8, 16\}& \{1, 7, 12\}& \{3, 11, 13\}& \{5, 9, 15\}& \{6, 10, $a$\}& \{14, 17, $b$\}\\
    & \{1, 15, 17\}& \{3, 5, 10\}& \{4, 16, $a$\}& \{6, 9, 14\}& \{7, 8, 11\}& \{12, 13, $b$\}\\
    & \{0, 9, 11\}& \{1, 13, 16\}& \{3, 6, $b$\}& \{4, 7, 14\}& \{5, 17, $a$\}& \{8, 12, 15\}\\
    & \{0, 1, $a$\}& \{4, 8, 9\}& \{7, 15, $b$\}& \{10, 12, 14\}& \{11, 16, 17\}\\
        & \{0, 13, 14\}& \{1, 6, 8\}& \{3, 4, 15\}& \{5, 12, 16\}& \{7, 9, 10\}\\%& \{11, $a$, $b$\}\\
    & \{0, 6, 15\}& \{1, 3, 9\}& \{4, 10, 17\}& \{5, 7, 13\}\\
\end{tabular}
\end{center}}
We list below the block sets of the derived designs $\mathrm{gcGDD}(2,3,18)$ of type $2^{9}$ at the point $a$ and $b$; the blocks in each row form a $\mathrm{PPC}$ of $I_{18}\setminus G$ for some $G\in\mathcal{G}$. Note that $\mathcal{B}_{a}$ and $\mathcal{B}_{b}$ have the same special group $\{0,9\}$.
{\small\begin{center}
\begin{tabular}{l l l l l l l l }
$\mathcal{B}_{a}$:
    & \{0, 13, 16\}& \{2, 8, 14\}& \{3, 4, 11\}& \{5, 6, 9\}& \{7, 15, 17\}\\%& \{1, 10, $b$\}
    & \{0, 5, 10\}& \{1, 16, 17\}& \{4, 9, 15\}& \{6, 7, 8\}& \{12, 13, 14\}\\%& \{2, 11, $b$\}
    & \{0, 7, 11\}& \{1, 4, 6\}& \{2, 5, 17\}& \{8, 9, 13\}& \{10, 14, 16\}\\%& \{3, 12, $b$\}
    & \{0, 6, 12\}& \{1, 14, 15\}& \{2, 3, 7\}& \{5, 8, 11\}& \{9, 10, 17\}\\%& \{4, 13, $b$\}
    & \{0, 1, 2\}& \{3, 9, 16\}& \{4, 8, 10\}& \{6, 13, 17\}& \{11, 12, 15\}\\%& \{5, 14, $b$\}
    & \{0, 3, 17\}& \{1, 5, 12\}& \{2, 4, 16\}& \{7, 9, 14\}& \{10, 11, 13\}\\%& \{6, 15, $b$\}
    & \{0, 8, 15\}& \{1, 9, 11\}& \{2, 6, 10\}& \{3, 5, 13\}& \{4, 12, 17\}\\%& \{7, 16, $b$\}
    & \{0, 4, 14\}& \{1, 7, 13\}& \{2, 9, 12\}& \{3, 10, 15\}& \{6, 11, 16\}\\%& \{8, 17, $b$\}
    & \{2, 13, 15\}& \{3, 6, 14\}& \{4, 5, 7\}& \{8, 12, 16\}\\%& \{0, 9, $b$\}
    & \{1, 3, 8\}& \{5, 15, 16\}& \{7, 10, 12\}& \{11, 14, 17\}\\
\end{tabular}
\end{center}

\begin{center}
\begin{tabular}{l l l l l l l l }
$\mathcal{B}_{b}$:
    & \{0, 14, 15\}& \{2, 12, 13\}& \{4, 7, 9\}& \{5, 11, 17\}& \{6, 8, 16\}\\ %& \{1, 10, $a$\}
    & \{0, 6, 13\}& \{1, 9, 14\}& \{3, 4, 5\}& \{7, 8, 10\}& \{15, 16, 17\}\\ %& \{2, 11, $a$\}
    & \{0, 4, 17\}& \{1, 5, 7\}& \{2, 9, 16\}& \{8, 11, 14\}& \{10, 13, 15\}\\%& \{3, 12, $a$\}
    & \{0, 1, 8\}& \{2, 14, 17\}& \{3, 9, 15\}& \{5, 6, 10\}& \{11, 12, 16\}\\%& \{4, 13, $a$\}
    & \{0, 7, 12\}& \{1, 13, 17\}& \{2, 3, 6\}& \{4, 8, 15\}& \{9, 10, 11\}\\%& \{5, 14, $a$\}
    & \{0, 5, 16\}& \{1, 2, 4\}& \{3, 10, 14\}& \{7, 11, 13\}& \{8, 9, 12\}\\%& \{6, 15, $a$\}
    & \{0, 2, 10\}& \{1, 11, 15\}& \{3, 8, 13\}& \{4, 12, 14\}& \{6, 9, 17\}\\%& \{7, 16, $a$\}
    & \{0, 3, 11\}& \{1, 6, 12\}& \{2, 7, 15\}& \{4, 10, 16\}& \{5, 9, 13\}\\%& \{8, 17, $a$\}
    & \{3, 7, 17\}& \{4, 6, 11\}& \{5, 12, 15\}& \{13, 14, 16\}\\%& \{0, 9, $a$\}
    & \{1, 3, 16\}& \{2, 5, 8\}& \{6, 7, 14\}& \{10, 12, 17\}\\
\end{tabular}
\end{center}}
\qed




{\Lemma\label{Dcqs20_2} There exists a $\mathrm{^{2}cDCQS}(2^{9}:2)$.}
\proof Let $X=I_{18}\cup\{a,b\}$ and $\mathcal{G} = \{\{i,i+9\}: 0\leq i \leq 8\}$.
We construct a $\mathrm{CQS}(2^{9}:2)$ on point set $X$ with stem $S=\{a,b\}$ and group set $\mathcal{G}$.
Let $\mathcal{H}$ be a permutation group generated by the following permutation:
$$ (0,3,6,9,12,15)(1,4,7,10,13,16)(2,5,8,11,14,17)(a)(b).$$
Our $\mathrm{CQS}(2^{9}:2)$ $(X,S,\mathcal{G},\mathcal{B})$ admits the automorphism group $\mathcal{H}$ and the base blocks are listed as follows. (The underlined blocks belong to short orbits of length $2$.)
\begin{center}
\begin{tabular}{l l l l l l l l }
\{0, 1, 2, $a$\}& \{0, 1, 3, 6\}& \{0, 1, 4, 5\}& \{0, 1, 7, 9\}& \{0, 1, 8, $b$\}\\
\{0, 1, 10, 15\}& \{0, 1, 11, 14\}& \{0, 1, 12, 17\}& \{0, 1, 13, 16\}& \{0, 2, 3, 5\}\\
\{0, 2, 4, 8\}& \{0, 2, 6, 11\}& \{0, 2, 7, 16\}& \{0, 2, 9, 12\}& \{0, 2, 10, $b$\}\\
\{0, 2, 13, 14\}& \{0, 3, 7, $a$\}& \{0, 3, 8, 12\}& \{0, 3, 10, 17\}& \{0, 3, 14, $b$\}\\
\{0, 4, 6, $b$\}& \{0, 4, 7, 12\}& \{0, 4, 9, 17\}& \{0, 4, 10, 14\}& \{0, 4, 11, 16\}\\
\{0, 5, 8, 16\}& \{0, 5, 10, $a$\}& \{0, 5, 11, 13\}& \{0, 5, 17, $b$\}& {\underline{\{0, 6, 12, $a$\}}}\\
\{0, 6, 16, 17\}& \{0, 7, 8, 10\}& \{0, 7, 11, 17\}& \{0, 7, 13, $b$\}& \{0, 8, 11, $a$\}\\
\{0, 8, 14, 17\}&  \{0, 13, 17, $a$\}& \{0, 14, 16, $a$\}& \{1, 2, 5, $b$\}&\{1, 2, 7, 10\}\\
\{1, 2, 11, 13\}& \{1, 2, 14, 17\}& \{1, 4, 8, 10\}& \{1, 4, 11, $a$\}&\{1, 4, 17, $b$\}\\
 \{1, 5, 8, 14\}& {\underline{\{1, 7, 13, $a$\}}}&{\underline{\{2, 8, 14, $a$\}}}\\
%\underline{\{2, 11, $a$, $b$\}}&\underline{\{0, 9, $a$, $b$\}}&\underline{\{1, 10, $a$, $b$\}}&
\end{tabular}
\end{center}%

The resolutions of derived designs at $0,1,2,a$ and $b$ are listed in Appendix C. Thus, we obtain a
$\mathrm{^{2}cDCQS}(2^{9}:2)$.
\qed
{\Lemma \label{RDGDD20} There exists an $\mathrm{RDGDD}(3,4, 10\{2\})$.}

\proof Let $X= I_{20}$ and $\mathcal{G} = \{\{i,i+10\}: 0\leq i \leq 9\}$. We construct a $\mathrm{GDD}(3,4, 20)$ of type $2^{10}$ on $X$ with group set $\mathcal{G}$. Let $\mathcal{H}$ be a permutation group generated by the following permutation:
$$(0,1,2,3,4,5,6,7,8,9)(10,11,12,13,14,15,16,17,18,19).$$
Our $\mathrm{GDD}(3,4,20)$ $(X,\mathcal{G},\mathcal{B})$ admits the automorphism group $\mathcal{H}$ and the base blocks are listed as follows:
\begin{center}
\begin{tabular}{l l l l l l l l }
\{0, 1, 2, 6\}&\{0, 1, 3, 4\}&\{0, 1, 12, 16\}&\{0, 1, 13, 14\}&\{0, 1, 15, 19\}\\
\{0, 1, 17, 18\}&\{0, 2, 4, 17\}&\{0, 2, 5, 18\}&\{0, 2, 7, 14\}&\{0, 2, 11, 16\}\\
\{0, 2, 13, 19\}&\{0, 3, 6, 18\}&\{0, 3, 11, 14\}&\{0, 3, 12, 19\}&\{0, 4, 11, 13\}\\
\{0, 4, 15, 16\}&\{0, 4, 18, 19\}&\{0, 5, 11, 19\}&\{0, 12, 14, 17\}&\{0, 12, 15, 18\}\\
\{0, 13, 15, 17\}&\{0, 13, 16, 18\}&\{10, 11, 12, 16\}&\{10, 11, 13, 14\}\\
\end{tabular}
\end{center}
For any $x\in G$, $G\in \mathcal{G}$, we can check that the derived design at the point $x$ of the $\mathrm{GDD}$ is an $\mathrm{RGDD}(2,3, 18)$ of type $2^{9}$ with group set $\mathcal{G}\setminus \{G\}$. We only need to list the resolutions of the derived designs at $0$ and $10$ to prove; the blocks in each row form a $\mathrm{PC}$ of $X\setminus \{0,10\}$.
{\small\begin{center}
\begin{tabular}{l l l l l l l l }
$\mathcal{B}_{0}$: & \{1, 2, 6\}&\{3, 7, 15\}&\{4, 18, 19\}&\{5, 14, 16\}&\{8, 11, 17\}&\{9, 12, 13\}\\
          &\{1, 3, 4\}&\{2, 7, 14\}&\{5, 8, 12\}&\{6, 17, 19\}&\{9, 11, 15\}&\{13, 16, 18\}\\
          &\{1, 7, 8\}&\{2, 11, 16\}&\{3, 12, 19\}&\{4, 5, 6\}&\{9, 14, 18\}&\{13, 15, 17\}\\
          &\{1, 12, 16\}&\{2, 13, 19\}&\{3, 5, 17\}&\{4, 8, 9\}&\{6, 14, 15\}&\{7, 11, 18\}\\
          &\{1, 5, 9\}&\{2, 4, 17\}&\{3, 11, 14\}&\{6, 8, 13\}&\{7, 16, 19\}&\{12, 15, 18\}\\
          &\{1, 17, 18\}&\{2, 3, 9\}&\{4, 15, 16\}&\{5, 7, 13\}&\{6, 11, 12\}&\{8, 14, 19\}\\
          &\{1, 15, 19\}&\{2, 5, 18\}&\{3, 8, 16\}&\{4, 11, 13\}&\{6, 7, 9\}&\{12, 14, 17\}\\
          &\{1, 13, 14\}&\{2, 8, 15\}&\{3, 6, 18\}&\{4, 7, 12\}&\{5, 11, 19\}&\{9, 16, 17\}\\
\end{tabular}
\end{center}

\begin{center}
\begin{tabular}{l l l l l l l l }
$\mathcal{B}_{10}$: & \{1, 2, 16\}&\{3, 15, 17\}&\{4, 9, 18\}&\{5, 6, 14\}&\{7, 8, 11\}&\{12, 13, 19\}\\
          &\{1, 3, 14\}&\{2, 4, 7\}&\{5, 9, 11\}&\{6, 13, 18\}&\{8, 12, 15\}&\{16, 17, 19\}\\
          &\{1, 4, 13\}&\{2, 5, 8\}&\{3, 16, 18\}&\{6, 9, 17\}&\{7, 12, 14\}&\{11, 15, 19\}\\
          &\{1, 5, 19\}&\{2, 15, 18\}&\{3, 6, 8\}&\{4, 12, 17\}&\{7, 9, 16\}&\{11, 13, 14\}\\
          &\{1, 6, 12\}&\{2, 3, 19\}&\{4, 5, 16\}&\{7, 13, 15\}&\{8, 9, 14\}&\{11, 17, 18\}\\
          &\{1, 8, 17\}&\{2, 9, 13\}&\{3, 5, 7\}&\{4, 6, 15\}&\{11, 12, 16\}&\{14, 18, 19\}\\
          &\{1, 9, 15\}&\{2, 14, 17\}&\{3, 4, 11\}&\{5, 12, 18\}&\{6, 7, 19\}&\{8, 13, 16\}\\
          &\{1, 7, 18\}&\{2, 6, 11\}&\{3, 9, 12\}&\{4, 8, 19\}&\{5, 13, 17\}&\{14, 15, 16\}\\
\end{tabular}
\end{center}}
\qed

We record a useful result on $\mathrm{RDS}$s.

\begin{Lemma}[\label{RDS}{\cite{RDS2,RDS}}]
There exists an $\mathrm{RDS}(3,q+1,q^{n}+1)$ for any prime power $q$ and any positive integer $n$.
\end{Lemma}

{\Theorem\label{result_1} There exists an $\mathrm{mcDSQS}(2\cdot 9^{n}+2)$ for any positive integer $n$.}
\proof For any positive integer $n$, apply Construction \ref{C1} with an $\mathrm{RDS}(3,10,9^{n}+1)$, which exists by Lemma \ref{RDS}. Since a $\mathrm{^{1}cDCQS}(2^{9}:2)$ exists by Lemma \ref{Dcqs20_1}, a $\mathrm{^{2}cDCQS}(2^{9}:2)$ exists by Lemma \ref{Dcqs20_2} and an $\mathrm{RDGDD}(3,4,10\{2\})$ exists by Lemma \ref{RDGDD20}, we can obtain a  $\mathrm{^{1}cDCQS}(2^{9^{n}}: 2)$. An $\mathrm{RDSQS}(4)$ exists obviously. Hence, an $\mathrm{mcDSQS}(2\cdot 9^{n}+2)$ exists by Construction \ref{cqs-wrdsqs1}.
\qed



\begin{Theorem}\label{result_3}
\begin{itemize}
 \item [$(1)$] Let $n\in\{2^{2m+1}+2:m\geq 0\}$. Then $\frac{n}{2}$ is the smallest $q$ such that a diameter perfect $(n,\frac{1}{4}\tbinom{n}{3},6;4)_{q}$ code exists.
 \item [$(2)$] Let $n \in\{ 2\cdot 9^{m}+2:m\geq 1\}\cup\{26,32\}$. Then $\frac{n}{2}+1$ is the smallest $q$ such that a diameter perfect $(n,\frac{1}{4}\tbinom{n}{3},6;4)_{q}$ code exists.
 \end{itemize}
\end{Theorem}
\proof Combine Theorem \ref{result_2} and Corollary \ref{shi-c1} to obtain the desired result in (1). Combine Lemma \ref{MIN_Dsqs}, Theorem \ref{result_1} and Corollary \ref{shi-c1} to obtain the result in (2).
\qed

%\begin{Theorem}[\label{RDS_2}{\cite{LEI1}}]
%\begin{itemize}
% \item [$(1)$] There exists an $\mathrm{RDSQS}(3 \cdot 7^{n} + 1)$ and an $\mathrm{RDSQS}(3 \cdot 13^{n} + 1)$ for any nonnegative integer $n$.
%\item [$(2)$] There exists an $\mathrm{RDSQS}(7 \cdot 4^{n})$ for any positive integer $n$.
%\end{itemize}
%\end{Theorem}






\section{Summary and open problems}

In this paper, we considered the constructions and existence of $\mathrm{mcDSQS}(v)$s. There has been an amount of  work on RDSQS$(6n+4)$s,
 while the study on $\mathrm{mcDSQS}(6n+2)$s has just begun in this paper. We generalized a number of existing recursive constructions about $\mathrm{RDSQS}$s to produce
  general $\mathrm{mcDSQS}$s. %and displayed a number of recursive constructions via Steiner systems $\mathrm{S}(3,K,v)$ with certain properties.
In particular, a  construction for $\mathrm{gcDCQS}$s is developed, which is also new even for $\mathrm{RDCQS}$s, see Construction \ref{C4}.
%a new recursive construction applicable for $\mathrm{RDCQS}$s is also developed, see Construction \ref{C4}.
Moreover, we demonstrated several constructions aiming only at $\mathrm{mcDSQS}(6n+2)$s (in Subsection 2.2). As the main results, we obtained a new infinite family of $\mathrm{RDSQS}(6n+4)$s and displayed a new infinite family of $\mathrm{LKTS}$s. (More families of $\mathrm{RDSQS}$s and thus more $\mathrm{LKTS}$s could be produced by substituting known orders to the product construction \cite{Zhou-lkts,yuan_olkts}, but we did not study this in depth to avoid distraction.)
We obtained the first infinite family of $\mathrm{mcDSQS}(6n+2)$s and presented their applications in diameter perfect constant-weight codes. Because the chromatic index of an STS$(v)$ equals $(v+1)/2$ only when $v\equiv 1\pmod 6$ and $v\ge 19$, direct constructions for  auxiliary designs to build $\mathrm{mcDSQS}(6n+2)$s
usually require significant computation.
For future work, we pose the following problems.\vspace{-0.2cm}
\begin{itemize}
\item[(1)] Find more examples of $\mathrm{mcDSQS}(v)$s, for example, an mcDSQS$(38)$ and an RDSQS$(52)$ (the smallest unknown mcDSQS$(6n+2)$s and the smallest unknown
 RDSQS$(6n+4)$s, respectively).\vspace{-0.2cm}
\item[(2)] Find a $\mathrm{gcDCQS}(v^4:0)$ for $v=20,32$. Apply Corollary \ref{Rmul}. If a $\mathrm{gcDCQS}(20^4:0)$ exists, then there exists an $\mathrm{mcDSQS}(5\cdot 4^{n})$ for any integer $n\geq 3$. In the same way, the existence of a $\mathrm{gcDCQS}(32^4:0)$ will yield an $\mathrm{mcDSQS}(2^{2n+1})$ for any  integer $n\geq 4$.\vspace{-0.2cm}
\item[(3)] Construct a $\mathrm{gcDCQS}(18^{4}:2)$ and an $\mathrm{FDGDD}(3,4,5\{18\})$ or  $\mathrm{RDGDD}(3,4,5\{18\})$. By Construction \ref{C3} or \ref{C2} and Lemma \ref{RDS}, the existence of an $\mathrm{mcDSQS}(18\cdot 4^n+2)$ would be obtained by filling with an $\mathrm{mcDSQS}(20)$.\vspace{-0.2cm}
\item[(4)] Develop more effective recursive constructions for $\mathrm{mcDSQS}(v)$s.
\end{itemize}
\begin{thebibliography}{77}
\frenchspacing

\bibitem{DPD}
R. Ahlswede, H.K. Aydinian, L.H. Khachatrian, \emph{On perfect codes and related concepts}, Des. Codes
Cryptogr., \textbf{22} (2001), 221--237.

%\bibitem{STS19}
%C. J. Colbourn, A. D. Forbes, M. J. Grannell, T. S. Griggs, P. Kaski, P. R. J. $\ddot{\text{O}}$sterg$\mathring{\text{a}}$rd, D. A. Pike, O. Pottonen, \emph{Properties of the Steiner Triple Systems of Order $19$}, the electronic journal of combinatorics, \textbf{17} (2010), R98.

\bibitem{zheng}
Y. Chang, H. Zheng, J. Zhou, \emph{Existence of frame-derived H-designs}, Des. Codes
Cryptogr., \textbf{87} (2019), 1415--1431.

\bibitem{Zhou-lkts}
Y. Chang, J. Zhou, \emph{Large set of Kirkman triple systems and related designs}, J. Combin. Theory (A),
\textbf{120} (2013), 649--670.

\bibitem{HTS14}
C.J. Colbourn, \emph{Computing the chromatic index of Steiner triple systems}, The
Computer J., \textbf{25} (1982), 338--339.



\bibitem{Delsarte}
P. Delsarte, \emph{An algebraic approach to the association schemes of coding
theory}, Philips Res. Rep., Supplement,
\textbf{10} (1973), 1--97.

\bibitem{Etzion}
T.  Etzion, \emph{Non-binary diameter perfect constant-weight codes}, IEEE Trans. Inf. Theory, \textbf{68}
(2022), 891--904.

\bibitem{HSQS}
H. Hanani, \emph{On quadruple systems}, Can. J. Math., \textbf{12} (1960), 145--157.

\bibitem{sfan}
A. Hartman, \emph{The fundamental construction for $3$-designs}, Discrete Math., \textbf{124} (1994), 107--132.


\bibitem{J_TD}
L. Ji, J. Yin, \emph{Constructions of new orthogonal arrays and covering arrays of strength three}, J. Combin. Theory (A),
\textbf{117} (2010), 236--247.

\bibitem{TSTS}
T. P. Kirkman, \emph{On a problem in combinations}, Cambridge and Dublin Math. J., \textbf{2} (1847), 191--204.

\bibitem{LEI1}
Y. Liu, J. Lei, \emph{More results on large sets of Kirkman triple systems}, Des. Codes Cryptogr., \textbf{91} (2023), 2677--2686.

\bibitem{KTS1}
J. Lu, \emph{Collected Works of Lu Jiaxi on Combinatorial Designs}, Inner Mongolia People$^{\prime}$s Press, Huhhot, Mongolia, 1990.




\bibitem{Cande}
H. Moh$\acute{\text{a}}$csy, D.K. Ray-Chaudhuri, \emph{Candelabra systems and designs}, J. Statist. Plann. Inference,
\textbf{106} (2002), 419--448.

\bibitem{CSQS}
K.T. Phelps, \emph{On cyclic Steiner systems $\mathrm{S}(3, 4,20)$}, Ann. Discrete Math., \textbf{7} (1980), 277--300.

\bibitem{KTS2}
D.K.  Ray-Chaudhuri, R.M. Wilson, \emph{Solution of Kirkman's schoolgirl problem}, Proc. Sympos. Pure Math., \textbf{19} (1971), 187--204.

\bibitem{SHI}
M.  Shi, Y.  Xia, D.S.  Krotov,
\emph{A family of diameter perfect constant-weight codes from Steiner systems},
J. Combin. Theory (A),
\textbf{200} (2023), 105790.

\bibitem{KF}
D.R. Stinson, \emph{A survey of Kirkman triple systems and related designs}, Discrete Math., \textbf{92} (1991), 371--393.

\bibitem{frame}
L. Teirlinck, \emph{Some new $2$-resolvable Steiner quadruple systems}, Des. Codes Cryptogr., \textbf{4} (1994), 5--10.

\bibitem{HTS}
S.A. Vanstone, D.R. Stinson, P.J. Schellenberg, A. Rosa, R.S. Rees, C.J. Colbourn, M.W. Carter, J.E. Carter, \emph{Hanani triple systems}, Israel J. Math., \textbf{83} (1993), 305--319.

\bibitem{RDS2}
E. Witt, \emph{$\ddot{U}$ber steinersche systeme}, Abhandlungen Mathematischen Seminar Univ. Hamburg, \textbf{12} (1938), 265--275.



\bibitem{DRSQS11}
J.  Xu, J.  Bao, L.  Ji, \emph{Doubly resolvable Steiner quadruple systems of orders $2^{2n+1}$}, Des. Codes Cryptogr., \textbf{88} (2020), 2377--2386.


\bibitem{Kang-2008}
L. Yuan, Q. Kang, \emph{Another construction for large sets of Kirkman triple systems}, Des. Codes Cryptogr., \textbf{48} (2008), 35--42.

\bibitem{yuanotripling}
L. Yuan and  Q. Kang, {\it A tripling construction for overlarge
sets of KTS*,} Discrete Math., \textbf{309} (2009) 975--981.

\bibitem{yuan_olkts}
L. Yuan and  Q. Kang, {\it On overlarge sets of Kirkman triple systems}, Discrete Math., \textbf{310} (2010) 2119--2125.

\bibitem{RDS}
J.  Zhou, Y. Chang, \emph{A new result on Sylvester$^{\prime}$s problem}, Discrete Math., \textbf{331} (2014), 15--19.
\end{thebibliography}
%\clearpage


\appendix
\section{The derived designs in Lemma \ref{CQS}}
{\noindent For the given $\mathrm{CQS}(8^{4}:2)$ in Lemma \ref{CQS}, we list below the block set of the derived $\mathrm{KTS}(33,9)$ at the point $0$.
The first four parts form $\mathrm{PC}$s of $\mathbb{Z}_{32}\setminus G_0$; other parts form $\mathrm{PC}$s of $(\mathbb{Z}_{32}\cup S)\setminus \{0\}$.}
%For any $x\in G$, $G\in \mathcal{G}$, we can check that the derived design at $x$ of the $\mathrm{CQS}(8^{4}:2)$ is a $\mathrm{KTS}(33,9)$ with the hole $(\mathcal{G}\cup S)\setminus \{x\}$. We only need to list the resolutions of the derived design at $0$ to prove;
{\small\begin{center}
\begin{tabular}{l l l l l l l l }
$\mathcal{B}_{0}$: &\{1, 7, 9\}& \{2, 13, 22\}& \{3, 17, 18\}& \{5, 11, 26\}& \{6, 21, 27\}& \{10, 19, 30\}\\
\smallskip
& \{14, 15, 29\}& \{23, 25, 31\}\\
&\{1, 10, 25\}& \{2, 3, 19\}& \{5, 9, 14\}& \{6, 11, 17\}& \{7, 22, 31\}& \{13, 29, 30\}\\
\smallskip
& \{15, 21, 26\}& \{18, 23, 27\}\\
&\{1, 14, 19\}& \{2, 11, 27\}& \{3, 9, 26\}& \{5, 21, 30\}& \{6, 23, 29\}& \{7, 10, 15\}\\
\smallskip
& \{13, 18, 31\}& \{17, 22, 25\}\\
&\{1, 15, 18\}& \{2, 25, 26\}& \{3, 10, 27\}& \{5, 22, 29\}& \{6, 7, 30\}& \{9, 19, 21\}\\
\smallskip
& \{11, 13, 23\}& \{14, 17, 31\}\\
&\{1, 31, $b$\}& \{2, 4, 18\}& \{3, 21, 24\}& \{5, 20, 25\}& \{6, 9, 15\}& \{7, 12, 27\}\\
\smallskip
& \{8, 11, 29\}& \{10, 16, 22\}& \{13, 19, $a$\}& \{14, 28, 30\}& \{17, 23, 26\}\\
&\{1, 24, 26\}& \{2, 17, $b$\}& \{3, 11, 14\}& \{4, 6, 10\}& \{5, 23, 28\}& \{7, 13, 21\}\\
\smallskip
& \{8, 27, 30\}& \{9, 20, 22\}& \{12, 15, 19\}& \{16, 29, 31\}& \{18, 25, $a$\}\\
&\{1, 21, 22\}& \{2, 16, 30\}& \{3, 29, $b$\}& \{4, 7, 19\}& \{5, 8, 15\}& \{6, 18, 20\}\\
\smallskip
& \{9, 23, $a$\}& \{10, 11, 31\}& \{12, 14, 26\}& \{13, 25, 28\}& \{17, 24, 27\}\\
&\{1, 17, 30\}& \{2, 15, 31\}& \{3, 8, 25\}& \{4, 14, 22\}& \{5, 27, $b$\}& \{6, 16, 26\}\\
\smallskip
& \{7, 24, 29\}& \{9, 12, 13\}& \{10, 18, 28\}& \{11, 21, $a$\}& \{19, 20, 23\}\\
&\{1, 2, $a$\}& \{3, 5, 13\}& \{4, 17, 29\}& \{6, 19, $b$\}& \{7, 11, 28\}& \{8, 14, 21\}\\
\smallskip
& \{9, 24, 31\}& \{10, 20, 26\}& \{12, 18, 30\}& \{15, 22, 23\}& \{16, 25, 27\}\\
&\{1, 5, 6\}& \{2, 12, 21\}& \{3, 4, 23\}& \{7, 25, $b$\}& \{8, 10, 13\}& \{9, 28, 29\}\\
\smallskip
& \{11, 20, 30\}& \{14, 16, 18\}& \{15, 17, $a$\}& \{19, 22, 24\}& \{26, 27, 31\}\\
&\{1, 11, 12\}& \{2, 10, 29\}& \{3, 7, 20\}& \{4, 5, 31\}& \{6, 14, 25\}& \{8, 19, 26\}\\
\smallskip
& \{9, 18, $b$\}& \{13, 15, 16\}& \{17, 21, 28\}& \{22, 27, $a$\}& \{23, 24, 30\}\\
&\{1, 20, 29\}& \{2, 7, 23\}& \{3, 6, $a$\}& \{4, 26, 30\}& \{5, 12, 17\}& \{8, 18, 22\}\\
\smallskip
& \{9, 11, 16\}& \{10, 21, $b$\}& \{13, 14, 27\}& \{15, 24, 25\}& \{19, 28, 31\}\\
&\{1, 4, 13\}& \{2, 6, 28\}& \{3, 12, 31\}& \{5, 18, 19\}& \{7, 8, 17\}& \{9, 25, 30\}\\
\smallskip
& \{10, 14, 24\}& \{11, 22, $b$\}& \{15, 20, 27\}& \{16, 21, 23\}& \{26, 29, $a$\}\\
&\{1, 8, 23\}& \{2, 14, 20\}& \{3, 15, 28\}& \{4, 21, 25\}& \{5, 7, 16\}& \{6, 12, 22\}\\
\smallskip
& \{9, 10, 17\}& \{11, 18, 24\}& \{13, 26, $b$\}& \{19, 27, 29\}& \{30, 31, $a$\}\\
&\{1, 27, 28\}& \{2, 8, 9\}& \{3, 22, 30\}& \{4, 11, 15\}& \{5, 10, $a$\}& \{6, 13, 24\}\\
\smallskip
& \{7, 18, 26\}& \{12, 25, 29\}& \{14, 23, $b$\}& \{16, 17, 19\}& \{20, 21, 31\}\\
&\{1, 3, 16\}& \{2, 5, 24\}& \{4, 9, 27\}& \{6, 8, 31\}& \{7, 14, $a$\}& \{10, 12, 23\}\\
\smallskip
& \{11, 19, 25\}& \{13, 17, 20\}& \{15, 30, $b$\}& \{18, 21, 29\}& \{22, 26, 28\}\\
\end{tabular}
\end{center}
}
%For any $x\in S$, we can check that the derived design at $x$ of the $\mathrm{CQS}(8^{4}:2)$ is a $\mathrm{KF}(8^{4})$ with the group set $\mathcal{G}$.
We list below the block set of the derived $\mathrm{KF}(8^{4})$ at $a$; the blocks in each part form a $\mathrm{PC}$ of $\mathbb{Z}_{32}\setminus G$ for some $G\in \mathcal{G}$.
{\small\begin{center}
\begin{tabular}{l l l l l l l l }
$\mathcal{B}_{a}$: &\{1, 3, 18\}& \{2, 17, 19\}& \{5, 7, 22\}& \{6, 21, 23\}\\
\smallskip
& \{9, 11, 26\}& \{10, 25, 27\}& \{13, 15, 30\}& \{14, 29, 31\}\\
&\{1, 10, 19\}& \{2, 11, 25\}& \{3, 17, 26\}& \{5, 14, 23\}\\
\smallskip
& \{6, 15, 29\}& \{7, 21, 30\}& \{9, 18, 27\}& \{13, 22, 31\}\\
&\{1, 11, 22\}& \{2, 13, 23\}& \{3, 14, 25\}& \{5, 15, 26\}\\
\smallskip
& \{6, 17, 27\}& \{7, 18, 29\}& \{9, 19, 30\}& \{10, 21, 31\}\\
&\{1, 14, 27\}& \{2, 15, 21\}& \{3, 9, 22\}& \{5, 18, 31\}\\
\smallskip
& \{6, 19, 25\}& \{7, 13, 26\}& \{10, 23, 29\}& \{11, 17, 30\}\\
&\{0, 3, 6\}& \{2, 7, 12\}& \{4, 26, 31\}& \{8, 11, 14\}\\
\smallskip
& \{10, 15, 20\}& \{16, 19, 22\}& \{18, 23, 28\}& \{24, 27, 30\}\\
&\{0, 7, 14\}& \{2, 3, 4\}& \{6, 24, 31\}& \{8, 15, 22\}\\
\smallskip
& \{10, 11, 12\}& \{16, 23, 30\}& \{18, 19, 20\}& \{26, 27, 28\}\\
&\{0, 22, 27\}& \{2, 28, 31\}& \{3, 8, 30\}& \{4, 7, 10\}\\
\smallskip
& \{6, 11, 16\}& \{12, 15, 18\}& \{14, 19, 24\}& \{20, 23, 26\}\\
&\{0, 30, 31\}& \{2, 20, 27\}& \{3, 10, 28\}& \{4, 11, 18\}\\
\smallskip
& \{6, 7, 8\}& \{12, 19, 26\}& \{14, 15, 16\}& \{22, 23, 24\}\\
&\{0, 9, 23\}& \{1, 15, 24\}& \{3, 5, 20\}& \{4, 19, 21\}\\
\smallskip
& \{7, 16, 25\}& \{8, 17, 31\}& \{11, 13, 28\}& \{12, 27, 29\}\\
&\{0, 11, 21\}& \{1, 7, 20\}& \{3, 13, 24\}& \{4, 17, 23\}\\
\smallskip
& \{5, 16, 27\}& \{8, 19, 29\}& \{9, 15, 28\}& \{12, 25, 31\}\\
&\{0, 13, 19\}& \{1, 12, 23\}& \{3, 16, 29\}& \{4, 15, 25\}\\
\smallskip
& \{5, 11, 24\}& \{7, 17, 28\}& \{8, 21, 27\}& \{9, 20, 31\}\\
&\{0, 15, 17\}& \{1, 16, 31\}& \{3, 12, 21\}& \{4, 13, 27\}\\
\smallskip
& \{5, 19, 28\}& \{7, 9, 24\}& \{8, 23, 25\}& \{11, 20, 29\}\\
&\{0, 1, 2\}& \{4, 5, 6\}& \{8, 9, 10\}& \{12, 13, 14\}\\
\smallskip
& \{16, 17, 18\}& \{20, 21, 22\}& \{24, 25, 26\}& \{28, 29, 30\}\\
&\{0, 5, 10\}& \{1, 4, 30\}& \{2, 24, 29\}& \{6, 9, 12\}\\
\smallskip
& \{8, 13, 18\}& \{14, 17, 20\}& \{16, 21, 26\}& \{22, 25, 28\}\\
&\{0, 18, 25\}& \{1, 8, 26\}& \{2, 9, 16\}& \{4, 22, 29\}\\
\smallskip
& \{5, 12, 30\}& \{6, 13, 20\}& \{10, 17, 24\}& \{14, 21, 28\}\\
&\{0, 26, 29\}& \{1, 6, 28\}& \{2, 5, 8\}& \{4, 9, 14\}\\
\smallskip
& \{10, 13, 16\}& \{12, 17, 22\}& \{18, 21, 24\}& \{20, 25, 30\}\\
\end{tabular}
\end{center}
}
%\clearpage

\section{An $\mathrm{mcDSQS}(26)$ and an $\mathrm{mcDSQS}(32)$ in Lemma \ref{MIN_Dsqs}}
{\noindent We first construct a cyclic $\mathrm{SQS}(26)$ on point set $\mathbb{Z}_{26}$. The base blocks are listed as follows. (The underlined blocks belong to short orbits of length $13$.)}
\begin{center}
\begin{tabular}{l l l l l l l l }
\{0, 1, 2, 4\}& \{0, 1, 5, 23\}& \{0, 1, 6, 22\}& \{0, 1, 7, 21\}& \{0, 1, 8, 20\}& \{0, 1, 9, 15\}\\
\{0, 1, 10, 12\}& \{0, 1, 11, 19\}& \underline{\{0, 1, 13, 14\}}& \{0, 1, 16, 18\}& \{0, 1, 17, 24\}& \{0, 2, 5, 7\}\\
\{0, 2, 6, 9\}& \{0, 2, 8, 13\}& \{0, 2, 12, 22\}& \{0, 2, 14, 18\}& \{0, 2, 15, 20\}& \{0, 3, 6, 10\}\\
\{0, 3, 9, 18\}& \{0, 3, 11, 15\}& \{0, 3, 12, 17\}& \underline{\{0, 3, 13, 16\}}& \{0, 3, 14, 21\}& \{0, 4, 9, 19\}\\
\{0, 4, 11, 21\}& \underline{\{0, 4, 13, 17\}}& \underline{\{0, 6, 13, 19\}}\\
\end{tabular}
\end{center}
The above cyclic $\mathrm{SQS}(26)$ is an $\mathrm{mcDSQS}(26)$.
We list below the block set of the derived design at the point $0$; the blocks in each part form a $\mathrm{PPC}$ of $\mathbb{Z}_{26}\setminus\{0\}$.
{\small\begin{center}
\begin{tabular}{l l l l l l l l l l}
\{1, 2, 4\}&\{3, 5, 24\}&\{6, 7, 14\}&\{8, 10, 22\}\\
\smallskip
\{9, 17, 20\}&\{12, 18, 19\}&\{13, 15, 21\}&\{16, 23, 25\}\\
\{1, 3, 25\}&\{2, 5, 7\}&\{4, 13, 17\}&\{6, 15, 23\}\\
\smallskip
\{8, 9, 24\}&\{10, 14, 16\}&\{11, 12, 20\}&\{18, 21, 22\}\\
\{1, 5, 23\}&\{2, 3, 19\}&\{4, 6, 16\}&\{7, 8, 18\}\\
\smallskip
\{9, 12, 21\}&\{10, 20, 24\}&\{11, 14, 22\}&\{15, 17, 25\}\\
\{1, 6, 22\}&\{2, 10, 11\}&\{3, 12, 17\}&\{4, 15, 18\}\\
\smallskip
\{5, 9, 16\}&\{7, 13, 20\}&\{8, 14, 25\}&\{19, 21, 24\}\\
\{1, 7, 21\}&\{2, 14, 18\}&\{3, 20, 22\}&\{4, 5, 10\}\\
\smallskip
\{6, 11, 24\}&\{8, 15, 16\}&\{12, 13, 25\}&\{17, 19, 23\}\\
\{1, 8, 20\}&\{2, 16, 17\}&\{3, 11, 15\}&\{4, 12, 14\}\\
\smallskip
\{5, 21, 25\}&\{6, 13, 19\}&\{7, 9, 10\}&\{22, 23, 24\}\\
\{1, 9, 15\}&\{2, 21, 23\}&\{3, 13, 16\}&\{4, 7, 24\}\\
\smallskip
\{5, 6, 12\}&\{8, 11, 17\}&\{10, 18, 25\}&\{14, 19, 20\}\\
\{1, 10, 12\}&\{2, 24, 25\}&\{3, 14, 21\}&\{4, 20, 23\}\\
\smallskip
\{5, 8, 19\}&\{6, 17, 18\}&\{7, 11, 16\}&\{9, 13, 22\}\\
\{1, 11, 19\}&\{2, 15, 20\}&\{3, 9, 18\}&\{4, 22, 25\}\\
\smallskip
\{5, 14, 17\}&\{6, 8, 21\}&\{10, 13, 23\}&\{12, 16, 24\}\\
\{1, 16, 18\}&\{2, 12, 22\}&\{3, 7, 23\}&\{4, 9, 19\}\\
\smallskip
\{5, 11, 13\}&\{6, 20, 25\}&\{10, 17, 21\}&\{14, 15, 24\}\\
\{1, 17, 24\}&\{2, 8, 13\}&\{3, 6, 10\}&\{4, 11, 21\}\\
\smallskip
\{5, 18, 20\}&\{7, 12, 15\}&\{9, 14, 23\}&\{16, 19, 22\}\\
\{1, 13, 14\}&\{2, 6, 9\}&\{3, 4, 8\}&\{5, 15, 22\}\\
\smallskip
\{7, 19, 25\}&\{11, 18, 23\}\\
\{7, 17, 22\}&\{8, 12, 23\}&\{9, 11, 25\}&\{10, 15, 19\}\\
\smallskip
\{13, 18, 24\}&\{16, 20, 21\}\\
\end{tabular}
\end{center}
}

%\section{Appendix: An $\mathrm{mcDCQS}(32)$ in Lemma \ref{MIN_Dsqs}}
Now, we construct a cyclic $\mathrm{SQS}(32)$ on point set $\mathbb{Z}_{32}$. The base blocks are listed as follows. (The underlined blocks belong to short orbits of length $16$ while the double-underlined blocks lie in short orbits of length $8$.)
\begin{center}
\begin{tabular}{l l l l l l l l }
\{0, 1, 2, 25\}& \{0, 1, 3, 29\}& \{0, 1, 4, 8\}& \{0, 1, 5, 27\}& \{0, 1, 6, 9\}& \{0, 1, 7, 30\}\\
\{0, 1, 10, 13\}& \{0, 1, 11, 15\}& \{0, 1, 12, 21\}& \{0, 1, 14, 17\}& \{0, 1, 16, 22\}& \{0, 1, 18, 26\}\\
\{0, 1, 19, 23\}& \{0, 1, 20, 28\}& \{0, 2, 4, 15\}& \{0, 2, 5, 26\}& \{0, 2, 6, 14\}& \{0, 2, 7, 23\}\\
\{0, 2, 8, 29\}& \{0, 2, 10, 21\}& \{0, 2, 11, 19\}& \{0, 2, 12, 20\}& \underline{\{0, 2, 16, 18\}}& \{0, 2, 17, 22\}\\
\{0, 2, 24, 27\}& \{0, 3, 12, 17\}& \{0, 3, 13, 25\}& \{0, 3, 14, 20\}& \{0, 3, 15, 28\}& \{0, 3, 16, 21\}\\
\{0, 4, 9, 18\}& \{0, 4, 10, 16\}& \{0, 4, 11, 25\}& \{0, 4, 17, 23\}& \{0, 4, 20, 27\}& \{0, 5, 12, 18\}\\
\{0, 5, 13, 22\}& \{0, 6, 13, 23\}& \{0, 7, 14, 22\}&\underline{\underline{\{0, 8, 16, 24\}}}
\end{tabular}
\end{center}
The above cyclic $\mathrm{SQS}(32)$ is an $\mathrm{mcDSQS}(32)$. We list below the block set of the derived design at the point $0$; the blocks in each part form a $\mathrm{PPC}$ of $\mathbb{Z}_{32}\setminus\{0\}$.
{\small\begin{center}
\begin{tabular}{l l l l l l l l l l}
\{2, 3, 9\}& \{4, 7, 19\}& \{5, 12, 18\}& \{6, 16, 17\}& \{8, 20, 22\}\\\smallskip
\{10, 14, 31\}& \{11, 28, 30\}& \{13, 15, 24\}& \{21, 27, 29\}& \{23, 25, 26\}\\

\{1, 3, 29\}& \{4, 5, 24\}& \{6, 12, 28\}& \{7, 11, 18\}& \{8, 17, 27\}\\\smallskip
 \{9, 20, 21\}& \{10, 15, 23\}& \{13, 16, 31\}& \{14, 19, 26\}& \{22, 25, 30\}\\

\{1, 7, 30\}& \{2, 4, 15\}& \{5, 14, 28\}& \{6, 29, 31\}& \{8, 18, 25\}\\\smallskip
 \{9, 23, 27\}& \{10, 17, 24\}& \{11, 13, 21\}& \{12, 19, 22\}& \{16, 20, 26\}\\

\{1, 2, 25\}& \{3, 5, 11\}& \{6, 18, 21\}& \{7, 14, 22\}& \{8, 26, 28\}\\\smallskip
 \{9, 19, 24\}& \{10, 12, 27\}& \{13, 17, 20\}& \{15, 16, 29\}& \{23, 30, 31\}\\

\{1, 5, 27\}& \{2, 7, 23\}& \{3, 14, 20\}& \{4, 10, 16\}& \{8, 19, 30\}\\\smallskip
\{9, 13, 26\}& \{11, 22, 24\}& \{12, 25, 29\}& \{15, 21, 31\}& \{17, 18, 28\}\\

\{1, 19, 23\}& \{2, 28, 31\}& \{3, 16, 21\}& \{4, 12, 30\}& \{5, 9, 25\}\\\smallskip
\{6, 22, 26\}& \{7, 13, 27\}& \{8, 14, 15\}& \{11, 17, 29\}& \{18, 20, 24\}\\

\{1, 10, 13\}& \{2, 5, 26\}& \{3, 4, 6\}& \{7, 12, 16\}& \{8, 21, 23\}\\\smallskip
\{9, 17, 30\}& \{11, 14, 27\}& \{18, 22, 31\}& \{19, 20, 29\}& \{24, 25, 28\}\\

\{1, 18, 26\}& \{2, 6, 14\}& \{3, 13, 25\}& \{4, 21, 22\}& \{5, 7, 29\}\\
\smallskip
\{8, 16, 24\}& \{9, 10, 28\}& \{11, 12, 23\}& \{15, 20, 30\}& \{19, 27, 31\}\\

\{1, 14, 17\}& \{2, 10, 21\}& \{3, 18, 19\}& \{4, 20, 27\}& \{5, 13, 22\}\\\smallskip
\{6, 8, 11\}& \{7, 15, 25\}& \{9, 12, 31\}& \{16, 23, 28\}& \{26, 29, 30\}\\

\{1, 4, 8\}& \{2, 12, 20\}& \{3, 7, 31\}& \{5, 15, 17\}& \{6, 27, 30\}\\\smallskip
\{9, 11, 16\}& \{10, 22, 29\}& \{13, 19, 28\}& \{14, 18, 23\}& \{21, 24, 26\}\\

\{1, 16, 22\}& \{2, 8, 29\}& \{3, 15, 28\}& \{4, 9, 18\}& \{5, 21, 30\}\\\smallskip
\{6, 13, 23\}& \{7, 17, 26\}& \{10, 19, 25\}& \{11, 20, 31\}& \{12, 14, 24\}\\

\{1, 6, 9\}& \{2, 24, 27\}& \{3, 22, 23\}& \{4, 11, 25\}& \{5, 8, 31\}\\\smallskip
\{7, 10, 20\}& \{12, 15, 26\}& \{13, 18, 29\}& \{14, 16, 30\}& \{17, 19, 21\}\\

\{1, 12, 21\}& \{2, 16, 18\}& \{3, 24, 30\}& \{4, 13, 14\}& \{5, 20, 23\}\\\smallskip
\{6, 15, 19\}& \{7, 8, 9\}& \{10, 11, 26\}& \{17, 25, 31\}& \{22, 27, 28\}\\

\{1, 24, 31\}& \{2, 13, 30\}& \{3, 8, 10\}& \{4, 17, 23\}& \{5, 16, 19\}\\\smallskip
\{6, 20, 25\}& \{7, 21, 28\}& \{9, 14, 29\}& \{15, 18, 27\}\\

\{1, 11, 15\}& \{2, 17, 22\}& \{3, 26, 27\}& \{4, 28, 29\}& \{6, 7, 24\}\\\smallskip
\{8, 12, 13\}& \{10, 18, 30\}& \{14, 21, 25\}\\

\{1, 20, 28\}& \{2, 11, 19\}& \{3, 12, 17\}& \{4, 26, 31\}& \{5, 6, 10\}\\\smallskip
\{9, 15, 22\}& \{16, 25, 27\}& \{23, 24, 29\}\\
\end{tabular}
\end{center}
}

\section{The derived designs in Lemma \ref{Dcqs20_2}}
{\noindent For the given $\mathrm{CQS}(2^{9}:2)$ in Lemma \ref{Dcqs20_2}, we list below the block set of the derived design at the point $x\in\{0,1,2\}$; the blocks in each row form a $\mathrm{PPC}$ of $X\setminus\{x\}$ and the last two rows form $\mathrm{PPC}$s of $I_{18}\setminus G$, where $x\in G$, $G\in \mathcal{G}$.}
{\small\begin{center}
\begin{tabular}{l l l l l l l l }
$\mathcal{B}_{0}$: & \{2, 9, 12\}& \{3, 14, $b$\}& \{4, 15, $a$\}& \{5, 11, 13\}& \{6, 16, 17\}& \{7, 8, 10\}\\
& \{1, 4, 5\}& \{2, 7, 16\}& \{3, 8, 12\}& \{6, 9, 14\}& \{11, 15, $b$\}& \{13, 17, $a$\}\\
& \{1, 12, 17\}& \{2, 3, 5\}& \{4, 6, $b$\}& \{7, 14, 15\}& \{8, 11, $a$\}& \{9, 10, 16\}\\
& \{1, 8, $b$\}& \{2, 15, 17\}& \{3, 9, 11\}& \{4, 7, 12\}& \{6, 10, 13\}& \{14, 16, $a$\}\\
& \{1, 7, 9\}& \{2, 6, 11\}& \{3, 4, 13\}& \{5, 10, $a$\}& \{8, 14, 17\}& \{12, 16, $b$\}\\
& \{1, 13, 16\}& \{2, 10, $b$\}& \{3, 7, $a$\}& \{4, 9, 17\}& \{5, 12, 14\}& \{6, 8, 15\}\\
& \{1, 11, 14\}& \{2, 4, 8\}& \{3, 10, 17\}& \{5, 9, 15\}& \{6, 12, $a$\}& \{7, 13, $b$\}\\
& \{1, 2, $a$\}& \{3, 15, 16\}& \{5, 17, $b$\}& \{8, 9, 13\}& \{10, 11, 12\}\\
& \{1, 3, 6\}& \{4, 10, 14\}& \{5, 8, 16\}& \{7, 11, 17\}& \{12, 13, 15\}\\%& \{9, $a$, $b$\}\\
& \{1, 10, 15\}& \{2, 13, 14\}& \{4, 11, 16\}& \{5, 6, 7\}\\
\end{tabular}
\end{center}}

{\small\begin{center}
\begin{tabular}{l l l l l l l l }
$\mathcal{B}_{1}$: & \{0, 2, $a$\}& \{3, 5, 10\}& \{4, 13, 14\}& \{6, 12, 16\}& \{7, 8, 17\}& \{9, 11, $b$\}\\
& \{0, 3, 6\}& \{2, 5, $b$\}& \{4, 11, $a$\}& \{8, 13, 15\}& \{9, 16, 17\}& \{10, 12, 14\}\\
& \{0, 4, 5\}& \{2, 3, 9\}& \{6, 14, 15\}& \{7, 12, $b$\}& \{8, 16, $a$\}& \{10, 13, 17\}\\
& \{0, 7, 9\}& \{2, 15, 16\}& \{3, 17, $a$\}& \{4, 8, 10\}& \{5, 11, 12\}& \{6, 13, $b$\}\\
& \{0, 10, 15\}& \{2, 11, 13\}& \{3, 7, 14\}& \{4, 17, $b$\}& \{5, 6, $a$\}& \{8, 9, 12\}\\
& \{0, 8, $b$\}& \{2, 14, 17\}& \{4, 6, 7\}& \{5, 9, 13\}& \{10, 11, 16\}& \{12, 15, $a$\}\\
& \{2, 4, 12\}& \{3, 8, 11\}& \{5, 15, 17\}& \{6, 9, 10\}& \{7, 13, $a$\}& \{14, 16, $b$\}\\
& \{0, 13, 16\}& \{2, 7, 10\}& \{3, 15, $b$\}& \{6, 11, 17\}& \{9, 14, $a$\}\\
& \{0, 11, 14\}& \{2, 6, 8\}& \{3, 12, 13\}& \{4, 9, 15\}& \{5, 7, 16\}\\%& \{10, $a$, $b$\}\\
& \{0, 12, 17\}& \{3, 4, 16\}& \{5, 8, 14\}& \{7, 11, 15\}\\
\end{tabular}
\end{center}
\begin{center}
\begin{tabular}{l l l l l l l l }
$\mathcal{B}_{2}$:
& \{0, 1, $a$\}& \{3, 11, 17\}& \{4, 9, 14\}& \{5, 7, 8\}& \{6, 10, 16\}& \{12, 13, $b$\}\\
& \{0, 3, 5\}& \{4, 10, 11\}& \{6, 12, 15\}& \{7, 9, 13\}& \{8, 14, $a$\}& \{16, 17, $b$\}\\
& \{0, 6, 11\}& \{1, 5, $b$\}& \{3, 7, 12\}& \{4, 13, 16\}& \{8, 10, 15\}& \{9, 17, $a$\}\\
& \{0, 13, 14\}& \{1, 6, 8\}& \{3, 16, $a$\}& \{4, 7, $b$\}& \{9, 11, 15\}& \{10, 12, 17\}\\
& \{0, 7, 16\}& \{1, 3, 9\}& \{4, 5, 17\}& \{8, 11, 12\}& \{10, 13, $a$\}& \{14, 15, $b$\}\\
& \{0, 15, 17\}& \{1, 11, 13\}& \{3, 8, $b$\}& \{4, 6, $a$\}& \{5, 9, 10\}& \{12, 14, 16\}\\
& \{0, 4, 8\}& \{1, 14, 17\}& \{5, 11, 16\}& \{6, 9, $b$\}& \{7, 15, $a$\}\\
& \{0, 10, $b$\}& \{1, 15, 16\}& \{3, 6, 13\}& \{5, 12, $a$\}& \{7, 11, 14\}\\
& \{1, 4, 12\}& \{3, 10, 14\}& \{5, 13, 15\}& \{6, 7, 17\}& \{8, 9, 16\}\\%& \{11, $a$, $b$\}\\
& \{0, 9, 12\}& \{1, 7, 10\}& \{3, 4, 15\}& \{5, 6, 14\}& \{8, 13, 17\}\\
\end{tabular}
\end{center}}
We list below the block sets of the derived designs at $a$ and $b$; the blocks in each row form a $\mathrm{PC}$ of $I_{18}$.
{\small\begin{center}
\begin{tabular}{l l l l l l l l }
$\mathcal{B}_{a}$:& \{0, 1, 2\}& \{3, 4, 5\}& \{6, 7, 8\}& \{9, 10, 11\}& \{12, 13, 14\}& \{15, 16, 17\}\\
& \{0, 3, 7\}& \{1, 9, 14\}& \{2, 10, 13\}& \{4, 12, 17\}& \{5, 8, 15\}& \{6, 11, 16\}\\
& \{0, 4, 15\}& \{1, 8, 16\}& \{2, 5, 12\}& \{3, 11, 14\}& \{6, 9, 13\}& \{7, 10, 17\}\\
& \{0, 5, 10\}& \{1, 4, 11\}& \{2, 7, 15\}& \{3, 8, 13\}& \{6, 14, 17\}& \{9, 12, 16\}\\
& \{0, 6, 12\}& \{1, 3, 17\}& \{2, 8, 14\}& \{4, 10, 16\}& \{5, 7, 9\}& \{11, 13, 15\}\\
& \{0, 8, 11\}& \{1, 12, 15\}& \{2, 9, 17\}& \{3, 6, 10\}& \{4, 7, 14\}& \{5, 13, 16\}\\
& \{0, 13, 17\}& \{1, 5, 6\}& \{2, 3, 16\}& \{4, 8, 9\}& \{7, 11, 12\}& \{10, 14, 15\}\\
& \{0, 14, 16\}& \{1, 7, 13\}& \{2, 4, 6\}& \{3, 9, 15\}& \{5, 11, 17\}& \{8, 10, 12\}\\
\end{tabular}
\end{center}
\begin{center}
\begin{tabular}{l l l l l l l l }
$\mathcal{B}_{b}$:& \{0, 1, 8\}& \{2, 12, 13\}& \{3, 4, 11\}& \{5, 15, 16\}& \{6, 7, 14\}& \{9, 10, 17\}\\
& \{0, 2, 10\}& \{1, 9, 11\}& \{3, 5, 13\}& \{4, 12, 14\}& \{6, 8, 16\}& \{7, 15, 17\}\\
& \{0, 3, 14\}& \{1, 6, 13\}& \{2, 16, 17\}& \{4, 10, 15\}& \{5, 9, 12\}& \{7, 8, 11\}\\
& \{0, 4, 6\}& \{1, 14, 16\}& \{2, 3, 8\}& \{5, 7, 10\}& \{9, 13, 15\}& \{11, 12, 17\}\\
& \{0, 5, 17\}& \{1, 3, 15\}& \{2, 4, 7\}& \{6, 10, 12\}& \{8, 9, 14\}& \{11, 13, 16\}\\
& \{0, 7, 13\}& \{1, 2, 5\}& \{3, 6, 17\}& \{4, 9, 16\}& \{8, 12, 15\}& \{10, 11, 14\}\\
& \{0, 11, 15\}& \{1, 7, 12\}& \{2, 6, 9\}& \{3, 10, 16\}& \{4, 5, 8\}& \{13, 14, 17\}\\
& \{0, 12, 16\}& \{1, 4, 17\}& \{2, 14, 15\}& \{3, 7, 9\}& \{5, 6, 11\}& \{8, 10, 13\}\\
\end{tabular}
\end{center}}
\clearpage

\end{document}
